\BourbakiExercisesSection

\begin{enumerate}
\def\labelenumi{\arabic{enumi}.}
\item
  证明：若 E 是一个含无穷多个元素的域 K 上的向量空间 \(\neq\{0\}\) ，则 E 的生成系所成的集对于序关系 \(\supset\) 不是归纳的（作出 E 的生成系的一个递减序列 \((S_{n})\) ，使得诸 \(S_{n}\) 的交是空的）。
\item
  设 K 为一个域，L 为 K 的子域，I 为一个集合，且 \((x_{i})_{1\leqslant i\leqslant m}\) 为 \(K_{s}^{I}\) 中的一个自由向量族，使得每个 \(x_{i}\) 的坐标都属于 L。设 V 为 \(K_{s}^{I}\) 中由这些 \(x_{i}\) 生成的向量子空间；证明存在 I 的一个含 m 个元素的有限子集 J，使得对 V 的每个向量 \(z=(\zeta_{\alpha})_{\alpha\in\mathbb{I}}\) ，所有 \(\zeta_{\alpha}\) 都属于由 L 与 m 个元素 \(\zeta_{\beta}\) （其中 \(\beta\in J\) ）生成的 K 的子域。（通过将第 5 节的推论 3 应用于空间 \(L_{s}^{I}\) ，证明可化归于存在 m 个指标 \(\beta_{i}\in I\) （ \(1\leqslant i\leqslant m\) ）使得 \(\Pr_{\beta_{i}}(x_{j})=\delta_{ij}\) 的情形。）
\item
  \begin{enumerate}
  \def\labelenumii{(\alph{enumii})}
  \tightlist
  \item
    设 G 是一个群，M 是 G 的一个无限子集。证明：若 \(G'\) 是由 M 生成的 G 的子群，则 \(\text{Card}(G') = \text{Card}(M)\) 。
  \end{enumerate}
\end{enumerate}

\begin{enumerate}
\def\labelenumi{(\alph{enumi})}
\setcounter{enumi}{1}
\item
  设 \(\mathbf{K}\) 为一个域， \(\mathbf{M}\) 为 \(\mathbf{K}\) 的一个无限子集。证明：若 \(\mathbf{K}'\) 是 \(\mathbf{K}\) 的由 \(\mathbf{M}\) 生成的子域，则 \(\operatorname{Card}(\mathbf{K}') = \operatorname{Card}(\mathbf{M})\) 。（考虑子集的两个序列 \((\mathbf{A}_n)_{n \geqslant 0}, (\mathbf{P}_n)_{n \geqslant 0}\) （这些子集属于 \(\mathbf{K}\) ），使得 \(\mathbf{A}_0 = \mathbf{M}\) ， \(\mathbf{P}_n\) 是 \(\mathbf{K}^*\) 的由 \(\mathbf{A}_n \cap \mathbf{K}^*\) 生成的乘法子群，而 \(\mathbf{A}_{n+1}\) 是 \(\mathbf{K}\) 的由 \(\mathbf{P}_n\) 生成的加法子群；然后应用 (a)。）
\item
  设 K 为一个域，I 为一个无限集；证明 \(\operatorname{Card}(\mathbf{K}) \leqslant \dim(\mathbf{K}_{s}^{\mathbf{I}})\) 。（将其归结为 I = N 的情形，并用反证法论证。设 B 为 \(K_{s}^{N}\) 的满足 \(\operatorname{Card}(\mathbf{B}) < \operatorname{Card}(\mathbf{K})\) 的一组基，并设 L 为 K 中由 B 的所有元素的坐标生成的子域；则 \(\operatorname{Card}(\mathbf{L}) < \operatorname{Card}(\mathbf{K})\) 。构造 K 中元素的一个序列 \((\xi_{n})_{n \geqslant 1}\) ，使得对所有 n， \(\xi_{n}\) 不属于 K 中由 L 与诸 \(\xi_{i}\) （其指标 i \textless{} n ）生成的子域；然后应用习题 2，通过考虑点 \(x = (\xi_{n})\) （属于 \(K_{s}^{N}\) ）得出矛盾。）
\item
  由(c)可推出，对任意域K及任意无限集I，
\end{enumerate}

\[
\dim (\mathrm{K} _ {s} ^ {\mathrm{I}}) = (\operatorname{Card} (\mathrm{K})) ^ {\operatorname{Card} (\mathrm{I})}
\]

(``Erdös-Kaplansky 定理'').

4 给出一个无限序列 \((\mathbf{F}_n)\) 的例子，该序列由某个向量空间 \(\mathbf{E}\) 的向量子空间组成，使得 \(\operatorname{codim}\left(\bigcap_{n} \mathbf{F}_n\right) > \sum_{n} \operatorname{codim}(\mathbf{F}_n)\) 。（取 \(\operatorname{codim}(\mathbf{F}_n) = 1\) 对所有 \(n\) 成立，并利用习题 3 \((d)\) 。）

\begin{enumerate}
\def\labelenumi{\arabic{enumi}.}
\setcounter{enumi}{4}
\tightlist
\item
  设 \((\mathbf{H}_{\lambda})_{\lambda \in L}\) 是向量空间 \(\mathbf{E}\) （在域 \(\mathbf{K}\) 上）的一族超平面（经过 0），它构成 \(\mathbf{E}\) 的一个覆盖。
\end{enumerate}

\begin{enumerate}
\def\labelenumi{(\alph{enumi})}
\item
  证明：若 K 有限，则 \(\operatorname{Card}(\mathbf{L}) \geqslant 1 + \operatorname{Card}(\mathbf{K})\) ；若 K 无限，则 \(\operatorname{Card}(\mathbf{L}) \geqslant \aleph_{0} = \operatorname{Card}(\mathbf{N})\) 。（对 r 作归纳证明：若 K 至少有 r 个元素，则 E 不能是 r 个超平面的并。）用例子说明，在没有附加假设的情况下，这些不等式不能改进（当 K 无限时，考虑空间 \(\mathbf{E} = \mathbf{K}_{s}^{(\mathbf{N})}\) ）。
\item
  若 E 是有限维的，证明 \(\operatorname{Card}(\mathbf{L}) \geqslant 1 + \operatorname{Card}(\mathbf{K})\) 。（对 \(\dim(\mathbf{E})\) 作归纳论证。）这一结论在无限维时不再成立；考虑 \(K^{(N)}\)
\end{enumerate}

¶ 6. 设 S 是向量 K-空间 E 的一个非空有限子集，且不包含 0。假设存在一个整数 \(k \geqslant 1\) 具有如下性质：(*) 对所有 \(x \in S\) ，存在 \(S = \{x\}\) 到 \(h\) 个自由子集的一个划分，其中 \(h \leqslant k\) 。

对于这样的划分 \((\mathbf{N}_i)_{1\leqslant i\leqslant h}\) 与每个整数序列 \((r_j)_{1\leqslant j\leqslant m}\) （整数取自 \(\{1,h\}\) ），子空间 \(\mathrm{F}_{r_1r_2\dots r_m}\) （属于 \(\mathbf{E}\) ）按对 \(m\) 的归纳定义如下： \(\mathrm{F}_{r_1}\) 是由 \(\mathrm{N}_{r_1},\mathrm{F}_{r_1\dots r_p}\) 生成的子空间，后者为由 \(\mathrm{F}_{r_1\dots r_{p-1}}\) 与 \(\mathrm{N}_{r_p}\) 的交所生成的子空间。假设：

(**) 对所有 \(x \in S\) 以及每个划分 \((N_i)_{1 \leqslant i \leqslant h}\) （将 \(S = \{x\}\) 分为至多 \(k\) 个自由子集），至少存在一个序列 \((r_j)_{1 \leqslant j \leqslant m}\) ，使得

\[
x \notin \mathrm{F} _ {r _ {1} r _ {2} \dots r _ {m}}.
\]

\begin{enumerate}
\def\labelenumi{(\alph{enumi})}
\item
  设 \(n\) 是整数 \(m \geqslant 1\) （对 \(x \in S\) 的所有可能选取）中的最小者，并设 \((\mathbf{N}_i)_{1 \leqslant i \leqslant h}\) 和 \((r_j)_{1 \leqslant j \leqslant m}\) 满足条件 (*) 和 (**)。证明必然有 \(n = 1\) 。（用反证法论证，假设 \(x \notin F_{r_1 r_2 \ldots r_n}\) 且 \(n > 1\) 。则必有 \(x \in F_{r_n}\) 。为简单起见，记 \(E_p = F_{r_1 r_2 \ldots r_p}\) （ \(1 \leqslant p \leqslant n\) ），并写 \(x = \sum_i \alpha_i y_i\) ，其中 \(\alpha_i \in K\) ， \(y_i \in N_{r_n}\) ；至少有一个 \(y_i\) 不属于 \(E_{n-1}\) ，把它记作 \(y\) ；然后令 \(N = N_{r_n} - \{y\}\) ，令 \(N_i' = N_i\) （ \(1 \leqslant i \leqslant h\) 且 \(i \neq r_n\) ），令 \(N_{r_n}' = N \cup \{x\}\) ，使得 \((\mathbf{N}_i')_{1 \leqslant i \leqslant h}\) 是由自由子集组成的 \(S - \{y\}\) 的一个划分。定义 \(E_0' = E\) ，并对 \(1 \leqslant p \leqslant n\) 把 \(E_p'\) 定义为由 \(E_{p-1}'\) 与 \(N_{r_p}'\) 的交生成的子空间。现在 \(y \notin E_{n-1}\) ；令 \(q\) 为使 \(y \notin E_q\) 的最小整数；证明 \(E_q' \not\subset E_q\) 。令 \(s\) 为使 \(E_s' \not\subset E_s\) 的最小整数；证明 \(r_s = r_n\) ；最后由关系 \(y \in E_s, E_{s-1}' \subset E_{s-1}\) 与 \(E_s' \not\subset E_s\) 推出矛盾。
\item
  由 (a) 推出，在给定假设下，存在 S 到 h 个自由子集的一个划分，其中整数 \(h \leqslant k\) 。
\end{enumerate}

\begin{enumerate}
\def\labelenumi{\arabic{enumi}.}
\setcounter{enumi}{6}
\tightlist
\item
  设 S 是向量 K-空间 E 的一个非空有限子集，且不包含 0。假设存在一个整数 \(k \geqslant 1\) 使得对于每个子集
\end{enumerate}

S 的每个子集 T，有 \(\text{Card}(T) \leqslant k.\text{rg}(T)\) 。证明存在将 S 分为 h 个自由子集的一个划分，其中某个整数 \(h \leqslant k\) 。（证明 S 满足习题 6 的条件（ \(^{**}\) ），对 \(\text{Card}(S)\) 作归纳论证，并在所有子空间 \(F_{r_{1}\ldots r_{m}}\) 中取一个维数尽可能小的子空间；若 G 是这样的子空间，且 j 是使 \(\text{Card}(G \cap N_{j})\) 为最小的指标，则证明

\[
\begin{array}{r l} k. \operatorname{Card} (\mathrm{G} \cap \mathrm{N} _ {j}) & \leqslant \operatorname{Card} (\mathrm{G} \cap (\mathrm{S} - \{x \})) \\ & \leqslant \operatorname{Card} (\mathrm{G} \cap \mathrm{S}) \leqslant k. \operatorname{Card} (\mathrm{G} \cap \mathrm{N} _ {j}).) \end{array}
\]

\begin{enumerate}
\def\labelenumi{\arabic{enumi}.}
\setcounter{enumi}{7}
\item
  设 E 是一个向量空间。证明 E 的每个在 End(E) 中不是 0 的右因子的自同态都是满射的（§ 2，习题 7）。
\item
  设K是一个域。
\end{enumerate}

\begin{enumerate}
\def\labelenumi{(\alph{enumi})}
\item
  在向量空间 \(\mathbf{E} = \mathrm{K}_s^{(\mathbf{Z})}\) 中，设 \((e_n)_{n\in \mathbf{Z}}\) 为典范基；设 \(v\) 为 \(\mathbf{E}\) 的由关系 \(v(e_n) = e_{n + 1}\) （对所有 \(n\in \mathbf{Z}\) ）定义的自同构。若 \(u = 1_{\mathbb{E}} - v\) ，证明 \(u\) 是 \(\mathbf{E}\) 的一个单射自同态，但 \(\dim (\operatorname {Coker}(u)) = 1\) 。
\item
  在向量空间 \(\mathbf{E} = \mathrm{K}_s^{(\mathbb{N})}\) 中，设 \((e_n)_{n\in \mathbb{N}}\) 为典范基；设 \(v\) 为 \(\mathbf{E}\) 的由 \(v(e_0) = e_0\) ， \(v(e_n) = e_{n - 1} + e_n\) （对 \(n\geqslant 1\) ）定义的自同态。证明 \(v\) 是 \(\mathbf{E}\) 的一个自同构，且若 \(u = 1_{\mathbb{E}} - v\) ，则 \(u\) 是 \(\mathbf{E}\) 的满足 \(\dim (\mathrm{Ker}(u)) = 1\) 的一个满射自同态。
\item
  设 I 为一个非空集合。证明向量空间 \(E = K_{s}^{(I)}\) 的每个自同构 u 都可以写成两个自同构之差 u = v - w，除非 K 只有两个元素且 I 仅由一个元素组成。（注意只需证明某个特定的自同构 u 具有所述形式即可；当 K 是两个元素的域时，分别考虑 \(\text{Card}(I) = 2\) 与 \(\text{Card}(I) = 3\) 的情形。）
\end{enumerate}

设 E 是域 K 上的向量空间。证明 E 的每个自同态 u 要么是自同构，要么可以写成 u = v - w，其中 v 和 w 是 E 的自同构。（注意，u 总可以换成 \(s_{1} \circ u \circ s_{2}\) ，其中 \(s_{1}\) 和 \(s_{2}\) 是 E 的两个自同构。区分两种情形，视 \(\dim(\text{Ker}(u)) \leqslant \dim(\text{Coker}(u))\) 还是 \(\dim(\text{Ker}(u)) > \dim(\text{Coker}(u))\) 而定；当 \(\text{rg}(u)\) 有限时，第一种情形总成立。当

\[
\dim (\operatorname{Ker} (u)) \leqslant \dim (\operatorname{Coker} (u)),
\]

它可以简化为以下情形： \(\operatorname{Ker}(u) \cap \operatorname{Im}(u) = \{0\}\) ；若

\[
\dim (\operatorname{Ker} (u)) = \dim (\operatorname{Coker} (u)),
\]

还可以假设 \(\operatorname{Im}(u) + \operatorname{Ker}(u) = \mathbf{E}\) ，然后应用习题 9 (c)。若 \(\dim (\operatorname {Ker}(u)) < \dim (\operatorname {Coker}(u))\) ，则在 \(\mathbf{E}\) 中存在 \(\operatorname {Ker}(u)\) 的一个形如 \(\mathbf{W} \oplus \operatorname{Im}(u)\) 的补子空间，满足 \(\dim (\mathbf{W}) = \dim (\operatorname {Im}(u)) = \dim (\mathbf{E})\) ；应用开头的注记并取 \(s_1(u(\operatorname {Im}(u))) = \operatorname{Im}(u)\) ，并利用习题 9 (a) 和 (c)。当 \(\dim (\operatorname {Ker}(u)) > \dim (\operatorname {Coker}(u))\) 时，可化归于 \(\operatorname{Ker}(u) \subset \operatorname{Im}(u)\) 的情形；这次取 \(s_1(\operatorname{Ker}(u)) = u(\operatorname{Ker}(u))\) ，并利用前面诸情形及习题 9 的结果。）

\begin{enumerate}
\def\labelenumi{\arabic{enumi}.}
\setcounter{enumi}{10}
\tightlist
\item
  设 E、F 是域 K 上的两个向量空间， \(u: E \rightarrow F\) 是一个线性映射。证明对于 F 的每个向量子空间 V，
\end{enumerate}

\[
\dim (u ^ {- 1} (\mathbf {V})) = \dim (\mathbf {V} \cap \operatorname{Im} (u)) + \dim (\operatorname{Ker} (u)).
\]

\begin{enumerate}
\def\labelenumi{\arabic{enumi}.}
\setcounter{enumi}{11}
\tightlist
\item
  设 E、F 为两个向量空间，u、v 为从 E 到 F 的两个线性映射。
\end{enumerate}

\begin{enumerate}
\def\labelenumi{(\alph{enumi})}
\tightlist
\item
  证明以下不等式成立：
\end{enumerate}

\[
\begin{array}{c} \operatorname{rg} (v) \leqslant \operatorname{rg} (u + v) + \operatorname{rg} (u) \\ \operatorname{rg} (u + v) \leqslant \inf (\dim (\mathrm{E}), \dim (\mathrm{F}), \operatorname{rg} (u) + \operatorname{rg} (v)). \end{array}
\]

若E和F是有限维的，证明 \(\operatorname{rg}(u + v)\) 可以取满足不等式的任意整数值

\[
\left| \operatorname{rg} (u) - \operatorname{rg} (v) \right| \leqslant \operatorname{rg} (u + v) \leqslant \inf (\dim (\mathrm{E}), \dim (\mathrm{F}), \operatorname{rg} (u) + \operatorname{rg} (v)).
\]

\begin{enumerate}
\def\labelenumi{(\alph{enumi})}
\setcounter{enumi}{1}
\tightlist
\item
  证明
\end{enumerate}

\[
\dim (\operatorname{Ker} (u + v)) \leqslant \dim (\operatorname{Ker} (u) \cap \operatorname{Ker} (v)) + \dim (\operatorname{Im} (u) \cap \operatorname{Im} (v)).
\]

\begin{enumerate}
\def\labelenumi{(\alph{enumi})}
\setcounter{enumi}{2}
\tightlist
\item
  证明
\end{enumerate}

\[
\dim (\operatorname{Coker} (u + v)) \leqslant \dim (\operatorname{Coker} (u)) + \operatorname{rg} (v).
\]

（若 \(\mathbf{V}\) 是 \(\operatorname{Ker}(v)\) 在 \(\mathbf{E}\) 中的一个补子空间，注意

\[
u (\operatorname{Ker} (v)) \subset \operatorname{Im} (u + v) \quad \text { and } \quad \dim (u (\mathbf {V})) \leqslant \operatorname{rg} (v).)
\]

\begin{enumerate}
\def\labelenumi{\arabic{enumi}.}
\setcounter{enumi}{12}
\tightlist
\item
  设 E、F、G 为三个向量空间，且 \(u: \mathbf{E} \to \mathbf{F}\) , \(v: \mathbf{F} \to \mathbf{G}\) 为两个线性映射。
\end{enumerate}

\begin{enumerate}
\def\labelenumi{(\alph{enumi})}
\tightlist
\item
  证明存在将 E 分解为 \(\operatorname{Ker}(u)\) 与两个子空间 M, N 的直和，使得 \(\operatorname{Ker}(v \circ u) = \mathbf{M} \oplus \operatorname{Ker}(u)\) 和
\end{enumerate}

\[
\operatorname{Im} (v \circ u) = v (u (\mathbf {N})).
\]

\begin{enumerate}
\def\labelenumi{(\alph{enumi})}
\setcounter{enumi}{1}
\item
  \(\dim(\operatorname{Ker}(u)) \leqslant \dim(\operatorname{Ker}(v \circ u)) \leqslant \dim(\operatorname{Ker}(u)) + \dim(\operatorname{Ker}(v))\) ；若 \(\operatorname{Ker}(u)\) 和 \(\operatorname{Ker}(v)\) 是有限维的，证明 \(\dim(\operatorname{Ker}(v \circ u))\) 可以取满足上述不等式的每一个整数值。
\item
  以下等式成立：
\end{enumerate}

\[
\begin{array}{l} \operatorname{rg} (u) = \operatorname{rg} (v \circ u) + \dim (\operatorname{Im} (u) \cap \operatorname{Ker} (v)) \\ \operatorname{rg} (v) = \operatorname{rg} (v \circ u) + \operatorname{codim} _ {\mathbb {F}} (\operatorname{Im} (u) + \operatorname{Ker} (v)). \end{array}
\]

若 F 是有限维的，证明 \(\operatorname{rg}(v \circ u)\) 可以取满足该不等式的任意整数值

\[
\sup (0, \operatorname{rg} (u) + \operatorname{rg} (v) - n) \leqslant \operatorname{rg} (v \circ u) \leqslant \inf (\operatorname{rg} (u), \operatorname{rg} (v)).
\]

\begin{enumerate}
\def\labelenumi{\arabic{enumi}.}
\setcounter{enumi}{13}
\item
  设 E、F、G 为三个向量空间，且 \(u: \mathbf{E} \to \mathbf{F}\) , \(w: \mathbf{E} \to \mathbf{G}\) 两个线性映射。证明若 \(u(0) \subset w(0)\) ，存在一个线性映射 \(v\) ，使得 \(w = v \circ u\) （参见第2节，习题8）。
\item
  设E，F是域K上的两个向量空间， \(u: E \rightarrow F\) 是一个线性映射。
\end{enumerate}

\begin{enumerate}
\def\labelenumi{(\alph{enumi})}
\item
  下列条件等价：（1） \(\operatorname{Ker}(u)\) 是有限维的；（2）存在一个线性映射 \(v\) （由 \(\mathbf{F}\) 到 \(\mathbf{E}\) ），使得 \(v \circ u = 1_{\mathbf{E}} + w\) ，其中 \(w\) 是 \(\mathbf{E}\) 的一个有限秩自同态；（3）对每个向量空间 \(\mathbf{G}\) （在 \(\mathbf{K}\) 上）与每个线性映射 \(f: \mathbf{G} \to \mathbf{E}\) ，关系 \(\operatorname{rg}(u \circ f) < +\infty\) 蕴含 \(\operatorname{rg}(f) < +\infty\) 。
\item
  以下条件等价：（1）Coker(u)是有限维的；（2）存在一个线性映射v：F→E，使得 \(u \circ v = 1_{F} + w\) ，其中w是F的一个有限秩自同态；（3）对于K上的每个向量空间G和每个线性映射 \(g: F \to G\) ，关系 \(\operatorname{rg}(g \circ u) < +\infty\) 蕴含 \(\operatorname{rg}(g) < +\infty\) .
\end{enumerate}

\begin{enumerate}
\def\labelenumi{\arabic{enumi}.}
\setcounter{enumi}{15}
\tightlist
\item
  设 E、F 是 K 上的两个向量空间， \(u: E \to F\) 是一个线性映射。若 \(\operatorname{Ker}(u)\) 和 \(\operatorname{Coker}(u)\) 都是有限维的，则称 u 具有有限指标，且数 \(d(u) = \dim(\operatorname{Ker}(u)) - \dim(\operatorname{Coker}(u))\) 称为 u 的指标。
\end{enumerate}

设G是K上的第三个向量空间， \(v: F \rightarrow G\) 一个线性映射。证明：如果三个线性映射 u, v, \(v \circ u\) 中有两个具有有限指标，则第三个也具有有限指标，且 \(d(v \circ u) = d(u) + d(v)\) （利用习题13(a)）。

\begin{enumerate}
\def\labelenumi{\arabic{enumi}.}
\setcounter{enumi}{16}
\tightlist
\item
  设 E, F, G, H 是 K 上的四个向量空间，且 \(u: E \rightarrow F\) , \(v: F \rightarrow G\) , \(w: G \rightarrow H\) 三个线性映射。证明
\end{enumerate}

\[
\operatorname{rg} (v \circ u) + \operatorname{rg} (w \circ v) \leqslant \operatorname{rg} (v) + \operatorname{rg} (w \circ v \circ u).
\]

\begin{enumerate}
\def\labelenumi{\arabic{enumi}.}
\setcounter{enumi}{17}
\item
  设 E、F 是域 K 上的两个向量空间，u、v 是 E 到 F 的两个线性映射。要使存在 E 的自同构 f 和 F 的自同构 g 使得 \(v = g \circ u \circ f\) ，必要且充分的条件是 \(\operatorname{rg}(u) = \operatorname{rg}(v)\) , \(\dim(\operatorname{Ker}(u)) = \dim(\operatorname{Ker}(v))\) 与 \(\dim(\operatorname{Coker}(u)) = \dim(\operatorname{Coker}(v))\) 。
\item
  设 E 是域 K 上的一个向量空间。
\end{enumerate}

\begin{enumerate}
\def\labelenumi{(\alph{enumi})}
\item
  证明每个映射 \(f\) （由 \(\mathbf{E}\) 到 \(\mathbf{E}\) ）只要与 \(\mathbf{E}\) 的所有自同构可交换，就具有形式 \(x \mapsto \alpha x\) ，其中 \(\alpha \in \mathbf{K}\) （首先证明：对所有 \(x \in \mathbf{E}\) 存在 \(\rho(x) \in \mathbf{K}\) 使得 \(f(x) = \rho(x)x\) ，利用 \(f\) 与 \(\mathbf{E}\) 的每个使 \(x\) 不变的自同构可交换这一事实）。
\item
  设 \(g\) 是 \(\mathbf{E} \times \mathbf{E}\) 到 \(\mathbf{E}\) 的一个映射，使得对每个自同构 \(u\) （属于 \(\mathbf{E}\) ）都有 \(g(u(x), u(y)) = u(g(x, y))\) （对所有 \(x, y\) 属于 \(\mathbf{E}\) ）。证明存在两个元素 \(\alpha, \beta\) （属于 \(\mathbf{K}\) ），使得 \(g(x, y) = \alpha x + \beta y\) 在由线性无关元素组成的有序对 \((x, y)\) （这些元素属于 \(\mathbf{E}\) ）所成的集合上成立；此外，存在一个映射 \(\phi\) （由 \(\mathbf{K} \times \mathbf{K}\) 到 \(\mathbf{K}\) ），使得 \(g(\lambda x, \mu x) = \phi(\lambda, \mu)x\) 对所有 \(x \in \mathbf{E}\) 成立（同样的方法）。若进一步 \(g(u(x), u(y)) = u(g(x, y))\) 对每个自同态 \(u\) （属于 \(\mathbf{E}\) ）成立，则 \(g(x, y) = \alpha x + \beta y\) 对所有 \(x, y\) （属于 \(\mathbf{E}\) ）成立。推广到 \(\mathbf{E}^n\) 到 \(\mathbf{E}\) 的映射。
\end{enumerate}

\begin{enumerate}
\def\labelenumi{\arabic{enumi}.}
\setcounter{enumi}{19}
\tightlist
\item
  设 E 是一个向量空间。
\end{enumerate}

\begin{enumerate}
\def\labelenumi{(\alph{enumi})}
\item
  证明：若 M 和 N 是 E 的两个向量子空间，而 M'、N' 分别是 \(\mathbf{M}\) 和 \(\mathbf{N}\) 在 \(\mathbf{E}^*\) 中的正交，则 \(\mathbf{M} \cap \mathbf{N}\) 的正交是 \(\mathbf{M}' + \mathbf{N}'\) （参见 § 2，习题 9 (a)）。
\item
  若 E 是无限维的，证明存在超平面 \(H'\) （属于 \(E^{*}\) ），使得 E 中与 \(H'\) 正交的子空间化为 0（考虑一个包含对应于 E 的一组基的坐标形式的超平面）。
\item
  证明：若 \(\mathbf{E}\) 是无限维的，则存在子空间的一个无限族 \((\mathbf{V}_i)\) （这些子空间属于 \(\mathbf{E}\) ），使得若 \(\mathbf{V}_i'\) 是 \(\mathbf{E}^*\) 中与 \(\mathbf{V}_i\) 正交的子空间，则 \(\mathbf{E}^*\) 中与 \(\bigcap_{i} \mathbf{V}_i\) 正交的子空间不同于 \(\sum_{i} \mathbf{V}_i'\) 。
\item
  由 (b) 推出，若 E 是无限维的，则存在一个作为直和的分解 \(V' \oplus W'\) （这是 \(E^{*}\) 的一个分解），其中 \(W'\) 是有限维的，使得 E 的子空间 V、W 之和 \(V + W\) （其中 V、W 分别与 \(V'\) 和 \(W'\) 正交）不同于 E。
\item
  为使 E 的一个子空间是有限维的，必要且充分条件是它在 \(E^{*}\) 中的正交具有有限余维数。
\end{enumerate}

\begin{enumerate}
\def\labelenumi{\arabic{enumi}.}
\setcounter{enumi}{20}
\item
  用第 6 节定理 8 的记号，设 \(u\) 表示线性映射 \(x \mapsto (\langle x, x_i^* \rangle)\) （由 \(\mathbf{E}\) 到 \(\mathbf{K}_s^{\mathbf{I}}\) ），并设 \(y_0 = (\eta_i)\) 。 \(\mathbf{K}_d^{(\mathbf{I})}\) 与 \(\mathbf{K}_s^{\mathbf{I}}\) 的对偶的一个子空间在 \(\mathbf{K}_d^{(\mathbf{I})}\) 到其双对偶的典范映射下恒等；然后证明：若 \(\mathbf{N}'\) 是 \(^t u\) 的核，则定理 3（第 5 节）的条件表达了下述事实： \(y_0\) 与 \(\mathbf{N}'\) 和 \(\mathbf{K}_d^{(\mathbf{I})}\) 的交正交。当 \(\mathbf{I}\) 无限且方程组 (21)（第 5 节）具有有限秩时，证明这个交不同于 \(\mathbf{N}'\) （注意 \(\mathbf{N}'\) 这时具有有限余维数）。
\item
  设 E 和 F 是域 K 上的两个向量空间，u 是 E 到 F 的一个线性映射。若 V 是 E 的一个向量子空间，而 \(V'\) 是 V 在 \(E^{*}\) 中的正交，证明 \(u(V)\) 的对偶同构于 \({}^{t}u(F^{*})/(V'\cap{}^{t}u(F^{*}))\) 。若 \(W'\) 是 \(F^{*}\) 的一个使 \({}^{t}u(W')\) 为有限维的子空间，而 W 是 \(W'\) 在 F 中的正交，则 \({}^{t}u(W')\) 同构于空间 \(u(E)/(W\cap u(E))\) 的对偶。
\item
  证明：为使线性映射 \(u\) （由向量空间 \(E\) 到向量空间 \(F\) ）满足 \(\mathrm{rg}(t u) = \mathrm{rg}(u)\) ，必要且充分条件是 \(\mathrm{rg}(u)\) 有限（参见习题 3 (d)）。
\item
  设 E 是交换域 K 上有限维数为 n 的向量空间（ \textgreater{} 1）；证明，除非 n = 2 且 K 是两个元素的域，否则不存在仅依赖于 E 的向量空间结构的同构 \(\phi\) （从 E 到 \(E^{*}\) 上）。（注意，若 \(\phi\) 是这样一个同构，则必然对 E 中的 x、y 以及 E 的每个自同构 u 有 \(\langle x, \phi(y) \rangle = \langle u(x), \phi(u(y)) \rangle\) （集合论，IV，§ 1，第 5 节）。）
\item
  \begin{enumerate}
  \def\labelenumii{(\alph{enumii})}
  \tightlist
  \item
    设 K 为一个域，L 与 M 为两个集合；存在典范包含 \((\mathbf{K}_{s}^{\mathbf{L}})^{(\mathbf{M})}\subset(\mathbf{K}_{s}^{(\mathbf{M})})^{\mathbf{L}}\subset\mathbf{K}_{s}^{\mathbf{L}\times\mathbf{M}}\) 。证明：要使这些向量空间中的两个相等，必要且充分条件是集合 L、M 中有一个是有限的。
  \end{enumerate}
\end{enumerate}

\begin{enumerate}
\def\labelenumi{(\alph{enumi})}
\setcounter{enumi}{1}
\tightlist
\item
  设 \((\mathbf{E}_{\lambda})_{\lambda \in L}\) 是 \(K\) 上的一个右向量空间族，而 \((\mathbf{F}_{\mu})_{\mu \in M}\) 是 \(K\) 上的一个左向量空间族。为使典范映射 (26)（第 7 节）是双射，必要且充分条件是向量空间 \(\prod_{\lambda \in L} \mathbf{E}_{\lambda}, \prod_{\mu \in M} F_{\mu}\) 之一是有限维的（利用 (a)）。
\end{enumerate}

\begin{enumerate}
\def\labelenumi{\arabic{enumi}.}
\setcounter{enumi}{25}
\tightlist
\item
  \begin{enumerate}
  \def\labelenumii{(\alph{enumii})}
  \tightlist
  \item
    设 E、F 是域 K 上的两个左向量空间；为使典范同态 \(E^{*} \otimes_{K} F \to \operatorname{Hom}_{K}(E, F)\) 成为双射，必要且充分条件是空间 E、F 之一是有限维的。
  \end{enumerate}
\end{enumerate}

\begin{enumerate}
\def\labelenumi{(\alph{enumi})}
\setcounter{enumi}{1}
\tightlist
\item
  设 E 是一个右向量空间，F 是域 K 上的一个左向量空间。为使典范同态 \(E \otimes_{K} F \to \operatorname{Hom}_{K}(E^{*}, F)\) 成为双射，必要且充分条件是 E 是有限维的（利用习题 3 (d)）。
\end{enumerate}

\begin{enumerate}
\def\labelenumi{\arabic{enumi}.}
\setcounter{enumi}{26}
\tightlist
\item
  \begin{enumerate}
  \def\labelenumii{(\alph{enumii})}
  \tightlist
  \item
    在命题 16 (i)（第 7 节）的假设下，为使典范同态 (27)（第 7 节）成为双射，必要且充分条件是 E 或 F 为有限维（注意 \(\mathrm{Hom}_{\mathbb{K}}(\mathbf{E},\mathbf{G})\otimes_{\mathbb{L}}\mathbf{F}\) 的一个元素在 (27)（第 7 节）下的像把 E 映到 \(\mathbf{G}\otimes_{\mathbb{L}}\mathbf{F}\) 的一个子空间，它包含在形如 \(\mathbf{G}\otimes_{\mathbb{L}}\mathbf{F}'\) 的子空间中，其中 \(\mathbf{F}'\) 是 \(\mathbf{F}\) 的一个有限维子空间）。
  \end{enumerate}
\end{enumerate}

\begin{enumerate}
\def\labelenumi{(\alph{enumi})}
\setcounter{enumi}{1}
\tightlist
\item
  在命题 16 (ii)（第 7 节）的假设下，为使典范同态 (28)（第 7 节）成为双射，必要且充分条件是有序对 \((\mathbf{E}_1, \mathbf{E}_2), (\mathbf{E}_1, \mathbf{F}_1), (\mathbf{E}_2, \mathbf{F}_2)\) 之一由有限维向量空间组成（利用 (a) 及 § 4 的习题 7）。
\end{enumerate}

\begin{enumerate}
\def\labelenumi{\arabic{enumi}.}
\setcounter{enumi}{27}
\item
  设 \(\rho: \mathbf{K} \to \mathbf{A}\) 是域 \(\mathbf{K}\) 到环 \(\mathbf{A}\) 的一个单射同态，并设 \(\mathbf{E}\) 是 \(\mathbf{K}\) 上的一个左向量空间。为使典范映射 \((\mathbf{E}^*)_{(\mathbf{A})} \to (\mathbf{E}_{(\mathbf{A})})^*\) 是双射，必要且充分条件是 \(\mathbf{E}\) 是有限维的，或者 \(\mathbf{A}\) （视为 \(\mathbf{K}\) 上的右向量空间）是有限维的（利用习题 25 ( \(b\) )）。当这两种情形之一成立时，存在典范双射 \((\mathbf{E}^{**})_{(\mathbf{A})} \to (\mathbf{E}_{(\mathbf{A})})^{**}\) 。
\item
  设 A 是一个整环，K 是它的分式域，E 是一个 A-模，而 T(E) 是它的挠模。
\end{enumerate}

\begin{enumerate}
\def\labelenumi{(\alph{enumi})}
\item
  证明 \(\mathbf{E}\) 上的每个线性形式都在 \(\mathrm{T}(\mathrm{E})\) 上为零，从而对偶空间 \(\mathbf{E}^*\) 和 \((\mathrm{E} / \mathrm{T}(\mathrm{E}))^*\) 典范同构。
\item
  证明典范映射 \((\mathbf{E}^{*})_{(\mathbf{K})}\to (\mathbf{E}_{(\mathbf{K})})^{*}\) 是单射，并且当 \(\mathbf{E}\) 有限生成时它是双射。
\item
  给出一个无挠 A-模 E 的例子，使得 \(E^{*} = \{0\}\) 而 \(E_{(K)} \neq \{0\}\) 。
\end{enumerate}

¶ 30. 设 A 是一个整环，E 和 F 是两个 A-模。证明若 x 是 E 的自由元素而 y 是 F 的自由元素，则 \(x \otimes y \neq 0\) （在 \(E \otimes_{A} F\) 中）。（先将其归结为 E 和 F 都无挠的情形，再归结为 E 和 F 都有限生成的情形，并利用习题 29 (b) 证明存在 \(E \times F\) 到 K 的一个 A-双线性映射 f，使得 \(f(x, y) \neq 0\) 。）

*31. 设 \(\mathbf{K}_0\) 为一个交换域，A 为多项式环 \(\mathrm{K}_0[\mathrm{X},\mathrm{Y}]\) （ \(\mathbf{K}_0\) 上两个未定元的多项式环，它是一个整环），而 E 为 A 的理想 \(\mathrm{AX} + \mathrm{AY}\) 。证明在张量积 \(\mathbf{E} \otimes_{\mathbf{A}} \mathbf{E}\) 中，元素 \(\mathbf{X} \otimes \mathbf{Y} - \mathbf{Y} \otimes \mathbf{X}\) 是 \(\neq 0\) 的，并且 \(\mathrm{XY}(\mathbf{X} \otimes \mathbf{Y} - \mathbf{Y} \otimes \mathbf{X}) = 0\) （考虑 \(\mathbf{E} \times \mathbf{E}\) 到商模 \(\mathbf{A}/\mathbf{E}\) 的 A-双线性映射）。*

\begin{enumerate}
\def\labelenumi{\arabic{enumi}.}
\setcounter{enumi}{31}
\item
  将第 10 节的结果推广到 A 为非交换环且具有左分式域 K 的情形（I，§ 9，习题 15）。（注意，若 \(\xi_{i}\) ( \(1 \leqslant i \leqslant n\) ) 是 K 的元素，则存在 A 中的 \(\alpha \neq 0\) ，使得所有 \(\alpha\xi_{i}\) 都属于 A。）
\item
  设 A 是一个整环。A-模 E 称为可除的，如果对所有 \(x \in E\) 以及 A 中所有 \(\alpha \neq 0\) ，存在 \(y \in E\) 使得 \(\alpha y = x\) 。
\end{enumerate}

\begin{enumerate}
\def\labelenumi{(\alph{enumi})}
\item
  设 E、F 是两个 A-模。证明：若 E 是可除的，则 \(E \otimes_{A} F\) 也是可除的。若 E 是可除的且 F 是无挠的，则 \(\operatorname{Hom}_{\mathbf{A}}(\mathbf{E}, \mathbf{F})\) 是无挠的。若 E 是一个挠 A-模而 F 是可除的，则 \(E \otimes_{A} F = \{0\}\) 。若 E 是一个挠 A-模而 F 是无挠的，则 \(\operatorname{Hom}_{\mathbf{A}}(\mathbf{E}, \mathbf{F}) = \{0\}\) 。
\item
  证明每个内射 A-模（§ 2，习题 11）都是可除的；反之，每个可除的无挠 A-模都是内射的（利用 § 2 习题 11 的判据 (δ)）。
\end{enumerate}

\begin{enumerate}
\def\labelenumi{\arabic{enumi}.}
\setcounter{enumi}{33}
\tightlist
\item
  设 A 是一个整环，P 是一个投射 A-模， \(P'\) 是 P 的一个投射子模，而 \(j: P' \to P\) 是典范单射。证明对每个无挠 A-模 E，同态 \(j \otimes 1: P' \otimes_{A} E \to P \otimes_{A} E\) 是单射（归结为 E 有限生成的情形，并将 E 嵌入一个自由 A-模）。
\end{enumerate}

¶ 35. (a) 设 K 为一个域，E 为一个 (K, K)-双模。假设 E 作为 K 上的左向量空间和右向量空间的维数都等于同一个有限数 n。证明存在 E 的 n 个元素的一个族 \((e_{i})\) ，它对于 E 上的两种向量空间结构中的每一种都是基。（注意，若 \((b_{j})_{1\leq j\leq m}\) 是 E 的一个由 m\textless n 个元素组成的族，它对于 E 上的两种向量空间结构中的每一种都是自由的，而 V 是由 \((b_{j})\) 生成的左向量子空间，W 是右向量子空间，则要么 \(V+W\neq E\) ，要么 \(V\cap CW\) 和 \(W\cap CV\) 都非空；在后一情形，若 \(y\in V\cap CW\) , \(z\in W\cap CV\) , 则 \(y+z\) 与 \(b_{j}\) 一起构成 \(m+1\) 个元素的一个族，它对于 E 上的两种向量空间结构中的每一种都是自由的。）

\begin{enumerate}
\def\labelenumi{(\alph{enumi})}
\setcounter{enumi}{1}
\item
  设 F 是 E 的一个子双模，其在 K 上的两种向量空间结构具有相同的维数 p \textless{} n.～证明存在 E 的元素的一个族 \((e_{i})_{1 \leqslant i \leqslant n}\) ，它对于 E 的每一种向量空间结构都是 E 的一组基，并且使得 \((e_{i})_{1 \leqslant i \leqslant p}\) 是 F 在其两种向量空间结构下的基（方法相同）。
\item
  设 \((b_{j})_{1\leqslant j\leqslant n - 1}\) 是 \(n - 1\) 个元素（属于 \(\mathbf{E}\) ）的一个族，它对于 \(\mathbf{E}\) 上的两种向量空间结构中的每一种都是自由的。设 \(\mathbf{V}\) （相应地，W）是 \((b_{j})\) 在 \(\mathbf{E}\) （视为左（相应地，右）向量空间）中生成的超平面。
\end{enumerate}

证明：若 \(\mathbf{V} \subset \mathbf{W}\) （相应地， \(\mathbf{W} \subset \mathbf{V}\) ），则必然 \(\mathbf{V} = \mathbf{W}\) （注意，若 \(a \notin \mathbf{W}\) ，则 \(\lambda \in \mathbf{K}\) 中使 \(\lambda a \in \mathbf{W}\) 的那些 λ 所成的集是一个左理想）。

设 \(\mathbf{K}_0\) 是一个交换域，而 \(\mathbf{K} = \mathbf{K}_0(\mathbf{X})\) 是 \(\mathbf{K}_0\) 上一个未定元的有理函数域（第四章）。在 K 上定义一个 (K, K)-双模结构如下：乘积 \(t.u\) （一个元素 \(u \in \mathbf{K}\) 与一个左算子 \(t \in \mathbf{K}\) 的乘积）是有理函数 \(t(\mathbf{X})u(\mathbf{X})\) ；乘积 \(u.t\) （ \(u\) 与一个右算子 \(t \in \mathbf{K}\) 的乘积）是有理函数 \(u(\mathbf{X})t(\mathbf{X}^2)\) ；这样在 K 上定义的左（相应地，右）向量 K-空间结构是 1 维的（相应地，2 维的）。由此推出 (K, K)-双模的一些例子，使得这种双模的两种向量空间结构的维数是任意整数。

\begin{enumerate}
\def\labelenumi{(\alph{enumi})}
\setcounter{enumi}{1}
\tightlist
\item
  从 (a) 推出一个 (K, K)-双模 E 的例子，其两个向量空间结构具有相同的维数，并且存在 E 的一个子双模 F，它的两个向量空间结构不具有相同的维数。*
\end{enumerate}

设 E 是一个左向量空间（有限维或无限维均可），而 \(A_{i}\) ( \(1 \leqslant i \leqslant n\) ) 是 E 的向量子空间。假设对点的每个序列 \((a_{i})_{1 \leqslant i \leqslant n}\) （这些点属于 E）使得对所有 i 有 \(a_{i} \in A_{i}\) ，由 \(a_{i}\) 生成的 E 的向量子空间都具有维数 \(\leqslant m\) （m 是一个整数 \(\leqslant n\) ）。证明存在 E 的一个维数为 \(h \leqslant m\) 的子空间 W，它包含 \(h + (n - m)\) 个子空间 \(A_{i}\) 。（对 n 作归纳论证。首先证明（对固定的 n）可以只考虑 m \textless{} n 且对所有 i 有 \(\dim(A_{i}) \leqslant m\) 的情形；还可假设对所有 i 有 \(A_{i} \neq \{0\}\) ，且这些 \(A_{i}\) 并不都是

维数为 1 的。然后（对固定的 \(n\) ）对 \(d = \sum_{i=1}^{m} \dim(A_i)\) 作归纳论证。例如，若 \(\dim(A_n) \geqslant 2\) ，考虑一个 1 维子空间 \(B_n\) （包含于 \(A_n\) ），并将归纳假设应用于 \(A_1, \ldots, A_{n-1}, B_n\) ；由此得出结论：存在一个数 \(k\) 使得 \(1 \leqslant k \leqslant m\) ，以及 E 的一个 \(k\) 维子空间 U，它包含 \(k\) 个 \(A_i\) （例如 \(A_1, \ldots, A_k\) ）；如有必要，把 \(k\) 换成一个整数 \(k'\) （满足 \(1 \leqslant k' \leqslant k\) ），并同时证明可以假设存在向量 \(b_i \in A_i\) （对 \(1 \leqslant i \leqslant k\) ）使得 \(b_1, \ldots, b_k\) 构成 U 的一组基；然后投影到 E 中 U 的一个补子空间上，并利用归纳假设。）

\begin{enumerate}
\def\labelenumi{(\alph{enumi})}
\setcounter{enumi}{1}
\tightlist
\item
  设 K 是一个有限域，E 是 K 上的一个 N 维向量空间。设 \(U_{i}\) ( \(1 \leqslant i \leqslant n\) ) 是 E 的向量子空间，使得对 \([1, n]\) 的每个子集 H，诸 \(U_{i}\) （其中 \(i \in H\) ）的交具有维数 \(\leqslant N - \text{Card}(H)\) 。证明这些 \(U_{i}\) 的并不能等于 E（利用对偶性并借助 (a) 论证）。
\end{enumerate}

\begin{enumerate}
\def\labelenumi{\arabic{enumi}.}
\setcounter{enumi}{37}
\item
  设 K 为特征 0 的交换域，E 为 K 上的有限维向量空间。设 \(p_{1}, \ldots, p_{m}\) 为 E 到 E 的向量子空间上的投影算子，使得 \(1_{E} = p_{1} + p_{2} + \cdots + p_{m}\) 。证明 E 是这些 \(p_i(\mathbf{E})\) 的直和，并且这些 \(p_i\) 两两可交换（取两边的迹）。
\item
  设 K 为一个域，L 为 K 的子域，且 \((\mathrm{K}_{\alpha})_{\alpha\in I}\) 为 K 的子域的一个左有向族，使得 \(L=\bigcap_{\alpha\in I}K_{\alpha}\) 。设 E 为 K 上的向量空间，且 \((a_{i})_{1\leqslant i\leqslant m}\) 为 E 的元素的一个有限族；证明若该族 \((a_{i})\) 在 L 上是自由的，则存在一个指标 \(\alpha\) 使得 \((a_{i})\) 在 \(K_{\alpha}\) 上是自由的。（对 m-r 作归纳论证，其中 r 是族 \((a_{i})\) 在 K 上的秩。）
\end{enumerate}

\BourbakiExercisesSection

\begin{enumerate}
\def\labelenumi{\arabic{enumi}.}
\tightlist
\item
  设 K 为一个域， \(K'\) 为 K 的一个子域，而 V 是 K 上的一个非零右向量空间。
\end{enumerate}

\begin{enumerate}
\def\labelenumi{(\alph{enumi})}
\item
  设 \(\mathbf{V}'\) 是一个 \(\mathbf{K}'\) -结构（在 \(\mathbf{V}\) 上）。证明 \(\mathbf{K}'\) 等于由下述 \(\mu \in \mathbf{K}\) 组成的集：对所有 \(x \in \mathbf{V}'\) , \(x\mu \in \mathbf{V}'\) 。
\item
  设 \(\Gamma'\) 为 K 的保持 K' 的元素不变的自同构群。证明：若不存在 K 的不属于 K' 且在所有自同构 \(\sigma \in \Gamma'\) 下不变的元素，则不存在 V 的不属于 V' 且在 V 的每个保持 V' 的所有元素不变的双射双态（相对于 K 的一个自同构）下不变的元素。
\item
  反之，设 E 为 V 的一个子集，G 为保持 E 的所有元素不变的 V 的双射双态群；对所有 \(u \in G\) ，设 \(\sigma_{u}\) 为 K 的相应自同构，并令 \(\Gamma\) 为 K 的自同构群，即 G 在同态 \(u \mapsto \sigma_{u}\) 下的像。假设 G 不包含 V 的不同于恒等自同构的自同构，并且不存在 V 的不属于 E 且在所有双态 \(u \in G\) 下不变的元素。证明：若 \(K'\) 是 K 的在所有自同构 \(\sigma \in \Gamma\) 下不变的元素所成的子域，则 E 是 V 上的一个 \(K'\) -结构。（首先证明 E 是 V 的一个加法子群且包含 V 的一组基；令 \(K''\) 为由下述元素 \(\mu \in K\) 组成的集：关系 \(x \in E\) 蕴含 \(x\mu \in E\) 。证明 \(K''\) 的元素在所有 \(\sigma \in \Gamma\) 下不变，并且 \(K''\) 是 K 的一个子域；由此推出 E 是 V 上的一个 \(K''\) -结构。）
\end{enumerate}

\begin{enumerate}
\def\labelenumi{\arabic{enumi}.}
\setcounter{enumi}{1}
\tightlist
\item
  设 K 为一个域， \(K'\) 为 K 的一个子域，V 是 K 上的一个右向量空间， \(V'\) 是 V 上的一个 \(K'\) -结构，而 \(R'\) 是对偶 \(V'*\) 在 V 的对偶 \(V*\) 中的典范像（第 4 节）。为使 \(R'\) 是 \(V*\) 的一个生成系，必要且充分条件是 V 是有限维的，或者 K 是一个有限维的右向量 \(K'\) -空间（利用 § 7 的习题 25）。
\end{enumerate}

¶ 3. 设 K 是一个域，L 是 K 的一个子域，令 \(K_{L}\) 表示把 K 视为 L 上的左向量空间而得到的域； \(E = \text{End}_{L}(K_{L})\) 典范地具有一个 (K, K)-双模结构。对偶 \((K_{L})^{*}\) 包含在 E 中，并且是 E 的一个 (K, L)-子双模。当 L 包含在 K 的中心内时，E 上的两个 L-向量空间结构——分别把 E 的两个 K-向量空间结构的标量域限制到 L 而得到——是相同的。

此后假设 \(\mathbf{L}\) 包含于 \(\mathbf{K}\) 的中心，且 \(\mathbf{K}_{\mathbf{L}}\) 的维数有限并等于 \(n\) 。

\begin{enumerate}
\def\labelenumi{(\alph{enumi})}
\item
  证明 \((\mathrm{K}_{\mathrm{L}})^{*}\) 在其 K 上的左向量空间结构下是 1 维的（用两种方式计算 \((\mathrm{K}_{\mathrm{L}})^{*}\) 在 L 上的维数）。
\item
  证明 \(\mathbf{E}\) 是一个 \(n\) 维左向量空间（在 \(\mathbf{K}\) 上）（同样的方法）。
\item
  设 F 是 K 上的一个有限维左向量空间，而 \(F_{[L]}\) 是 L 上的相应向量空间。若 \(u_{0} \neq 0\) 是 \(K_{L}\) 上的一个线性形式，证明映射 \(x' \mapsto u_{0} \circ x'\) （由 F 的对偶 \(F^{*}\) （把 F 视为 L 上的向量空间）映到对偶 \((\mathrm{F}_{[\mathrm{L}]})^{*}\) （即 \(F_{[L]}\) 的对偶）上）是双射的（利用 (a)）。
\end{enumerate}

¶ 4. 设 K 是其中心 Z 上有限秩的域，L 是 K 的包含 Z 的子域。

\begin{enumerate}
\def\labelenumi{(\alph{enumi})}
\item
  证明K关于L的左、右向量空间结构具有相同的维数。
\item
  证明习题 3 的性质 (a)、(b)、(c) 在这种情形也成立（若 \(u_0\) 是一个线性形式 \(\neq 0\) （在 \(\mathbf{L}_{\mathbf{Z}}\) 上），而 \(v_0\) 是一个线性形式 \(\neq 0\) （在 \(\mathbf{K}_{\mathbf{L}}\) 上），注意 \(\xi. (u_0 \circ v_0) = u_0 \circ (\xi v_0)\) （对 \(\xi \in \mathbf{K}\) ），并且 \(\xi. (u_0 \circ v_0)\) 遍历对偶 \((\mathbf{K}_{\mathbf{Z}})^*\) （当 \(\xi\) 遍历 \(\mathbf{K}\) 时），并利用习题 3 (c)）。
\end{enumerate}

\begin{enumerate}
\def\labelenumi{\arabic{enumi}.}
\setcounter{enumi}{4}
\tightlist
\item
  设 \(K_{0}\) 是域 K 的一个子域，使得 K 作为 \(K_{0}\) 上的右向量空间具有维数 2。设 E 是 K 上的一个左向量空间。
\end{enumerate}

\begin{enumerate}
\def\labelenumi{(\alph{enumi})}
\item
  设 \(E_{0}\) 是 E 的一个子集，它是 \(K_{0}\) 上的向量空间，并设 V 是 \(E_{0}\) 的最大向量子 K-空间；若 \(W_{0}\) 是一个向量子 \(K_{0}\) -空间（属于 \(E_{0}\) ，与 V 互补），证明由 \(W_{0}\) 生成的 E 的向量子 K-空间 W 满足 \(V \cap W = \{0\}\) （注意，若元素 \(\mu \in K\) 不属于 \(K_{0}\) ，则关系 \(\mu x \in E_{0}\) （对某个 \(x \in E_{0}\) ）蕴含 \(x \in V\) ）。
\item
  设 \(E_{0}^{\prime}\) 是 E 的另一个向量子 \(K_{0}\) -空间，而 \(V^{\prime}\) 是 \(E_{0}^{\prime}\) 的最大向量子 K-空间。为使存在 E 的一个把 \(E_{0}\) 映到 \(E_{0}^{\prime}\) 的 K-自同构，必要且充分条件是 V 和 \(V^{\prime}\) 关于 K 具有相同的维数，V 在 \(E_{0}\) 中与 \(V^{\prime}\) 在 \(E_{0}^{\prime}\) 中的余维数（关于 \(K_{0}\) ）相等，并且 \(E_{0}\) 与 \(E_{0}^{\prime}\) 在 E 中的余维数（关于 \(K_{0}\) ）相等（利用 (a)）。
\end{enumerate}

\BourbakiExercisesSection

*1. 在仿射空间 \(\mathbf{E}\) （在域 \(\mathbf{K}\) 上）中，四元组 \((a, b, c, d)\) （ \(\mathbf{E}\) 的点的四元组）称为平行四边形，如果 \(b - a = c - d\) ，这时 \((a, d, c, b)\) 也是一个平行四边形。证明若 \(\mathbf{K}\) 的特征 \(\neq 2\) ，则有序对 \((a, c)\) 和 \((b, d)\) 的中点相等；当 \(\mathbf{K}\) 的特征为 \(2?_{*}\) 。

*2. 设 E 是特征 \(\neq 2\) 的域上的一个仿射空间，而 \(a, b, c, d\) 是 E 的任意四个点。证明若 \(x, y, z, t\) 分别是有序对 \((a, b), (b, c), (c, d), (d, a)\) 的中点，则四元组 \((x, y, z, t)\) 是一个平行四边形（习题 1）。*

*3. 设 K 是特征 \(\neq 2\) 的交换域，E 是 K 上的一个仿射平面，而 \(a, b, c, d\) 是 E 的四个点，其中任意三点都不共线。设 \(\mathbf{D}_{xy}\) 表示过 E 的两个不同点 \(x, y\) 的直线。假设直线 \(\mathbf{D}_{ab}\) 和 \(\mathbf{D}_{cd}\) 有一个公共点 \(e\) ，且 \(\mathbf{D}_{ad}\) 和 \(\mathbf{D}_{bc}\) 有一个公共点 \(f\) 。证明三个有序对 \((a, c), (b, d), (e, f)\) 的中点位于同一条直线上。当 \(\mathbf{D}_{ab}\) 和 \(\mathbf{D}_{cd}\) 平行或 \(\mathbf{D}_{ad}\) 和 \(\mathbf{D}_{bc}\) 平行时，这一性质变成什么？考虑 K 是 3 个元素的域的情形。*

*4. 设 K 是特征异于 2 和 3 的域，E 是 K 上的一个仿射空间， \(a, b, c\) 是 E 的不在同一条直线上的三个点，而 \(a', b', c'\) 分别是有序对 \((b, c)\) , \((c, a)\) 和 \((a, b)\) 的中点。证明（用习题 3 的记号）直线 \(\mathbf{D}_{aa'}\) , \(\mathbf{D}_{bb'}\) 和 \(\mathbf{D}_{cc'}\) 都过三个点 \(a, b, c\) 的重心。当 K 的特征为 2 或特征为 3 时，这一性质变成什么？推广到 \(n\) 个仿射无关点的系统的情形。*

\begin{enumerate}
\def\labelenumi{\arabic{enumi}.}
\setcounter{enumi}{4}
\item
  设 E 是域 K 上的一个仿射空间，且 K 至少含三个元素。为使 E 的非空子集 V 是一个线性簇，必要且充分条件是：对 V 中每对 \((x,y)\) 不同的点 x 和 y，过 x 和 y 的直线 \(D_{xy}\) 整个包含在 V 中。若 K 恰有两个元素，则为使 V 是一个线性簇，必要且充分条件是 V 的任意三点的重心属于 V。
\item
  \begin{enumerate}
  \def\labelenumii{(\alph{enumii})}
  \tightlist
  \item
    设 E 是域 K 上的一个 \(\geqslant2\) 维仿射空间。为使 E 到自身的仿射映射 u 把 E 的每条直线变为平行直线，必要且充分条件是：与 u 相伴的线性映射 v 是一个位似 \(t\mapsto\gamma t\) （比例 \(\gamma\neq0\) 属于 K 的中心）。若 \(\gamma=1\) ，u 是一个平移；证明若 \(\gamma\neq1\) ，存在且仅存在一个点 \(a\in E\) 使得 \(u(a)=a\) 。若取 a 作为 E 的原点，则 u 在 E 上由此确定的向量空间结构下等同于一个中心位似；u 称为仿射空间 E 的以 a 为中心、以 \(\gamma\) 为比例的中心位似。
  \end{enumerate}
\end{enumerate}

\begin{enumerate}
\def\labelenumi{(\alph{enumi})}
\setcounter{enumi}{1}
\item
  设 \(u_1, u_2\) 是 \(\mathbf{E}\) 到 \(\mathbf{E}\) 的两个仿射映射，其中每一个要么是平移，要么是中心位似。证明 \(u_1 \circ u_2\) 是 \(\mathbf{E}\) 的一个平移或一个中心位似；若 \(u_1, u_2\) 和 \(u_1 \circ u_2\) 这三个都是中心位似，证明它们的中心位于同一条直线上。当 \(u_1\) 和 \(u_2\) 是中心位似而 \(u_1 \circ u_2\) 是平移时，能说些什么？
\item
  证明平移集与中心位似之集构成仿射群的一个正规子群 H，并且 H/T 同构于 K 的中心的乘法群；证明 H 仅当 H = T 时才可能是交换的，换言之，仅当 K 的中心恰含两个元素时。
\end{enumerate}

¶ 7. 设 E（相应地 E'）是域 K（至少含 3 个元素）上的一个有限维 \(n \geqslant 2\) 仿射空间（相应地，域 K' 上的仿射空间），并设 u 是 E 到 E' 的一个单射映射，它把 E 中同一条直线上任意三个点映到同一条直线上的三个点，并且使得由 \(u(E)\) 在 E' 中生成的仿射线性簇等于 E'。

\begin{enumerate}
\def\labelenumi{(\alph{enumi})}
\item
  证明 \(u\) 把 \(E\) 的每个仿射无关点系统映为一个仿射无关点系统（利用习题 5）。
\item
  设 \(D_{1}, D_{2}\) 是 E 中两条平行直线，而 \(D_{1}^{\prime}, D_{2}^{\prime}\) 是 \(E^{\prime}\) 中分别包含 \(u(\mathbf{D}_{1})\) 和 \(u(\mathbf{D}_{2})\) 的直线。证明 \(D_{1}^{\prime}\) 和 \(D_{2}^{\prime}\) 位于同一平面内；若进一步 u 是满射，证明 \(D_{1}^{\prime}\) 和 \(D_{2}^{\prime}\) 平行（在相反情形，证明 E 中会有 3 个不在同一直线上的点，它们在 u 下的像在同一直线上）（参见习题 17）。
\item
  此后假设：若 \(D_{1}\) 和 \(D_{2}\) 是 E 中的平行直线，则 \(E'\) 中分别包含 \(u(\mathbf{D}_{1})\) 和 \(u(\mathbf{D}_{2})\) 的直线是平行的。证明若取 E 中的原点 a 及原点 \(a' = u(a)\) （在 \(E'\) 中），则存在一个同构 \(\sigma\) （把 K 映到一个子域 \(K_{1}\) （属于 \(K'\) ）上），使得若把 E 视为 K 上的向量空间，且把 \(E'\) 作为 \(K_{1}, u\) 上的向量空间，则 u 是一个单射半线性映射（相对于 \(\sigma\) ，由 E 到 \(E'\) ）（§ 1，第 13 节）。（先考虑 n = 2 的情形；给定 E 的一个基 \((e_{1}, e_{2})\) ，证明对 K 的任意两个元素 \(\alpha, \beta\) ，点 \((\alpha + \beta)e_{1}\) 和 \((\alpha\beta)e_{1}\) （属于 E）可以从点 0、 \(e_{1}, e_{2}, \alpha e_{1}, \beta e_{1}\) 出发，通过作给定直线的平行线及给定直线的交点而构造出来；由此推出 \(u(\lambda e_{1}) = \lambda^{\sigma}u(e_{1})\) ，其中 \(\sigma\) 是 K 到 \(K'\) 的一个子域上的同构，然后证明也有 \(u(\lambda e_{2}) = \lambda^{\sigma}u(e_{2})\) （通过考虑连接点 \(\lambda e_{1}\) 和 \(\lambda e_{2}\) 的直线）。最后由此过渡到 n 为任意的情形，对 n 作归纳论证。）若 u 是双射，证明 \(K_{1} = K'\) 。
\item
  把 (c) 的结果推广到 K 是含 2 个元素的域的情形，同时假设 u 也把 E 的每个仿射无关点系统映为 E' 的仿射无关点系统。
\end{enumerate}

\begin{enumerate}
\def\labelenumi{\arabic{enumi}.}
\setcounter{enumi}{7}
\tightlist
\item
  设 E 是域 K 上的左仿射空间，T 为其平移空间。
\end{enumerate}

\begin{enumerate}
\def\labelenumi{(\alph{enumi})}
\tightlist
\item
  为使映射 \(f\) （由 \(\mathbf{E}\) 到一个左向量空间 \(\mathbf{L}\) （在 \(\mathbf{K}\) 上））是仿射的，必要且充分条件是
\end{enumerate}

\[
\begin{array}{r l} f (\mathbf {t} + x) - f (x) & = f (\mathbf {t} + y) - f (y) \\ f (\lambda \mathbf {t} + x) - f (x) & = \lambda f (\mathbf {t} + x) - f (x)) \end{array}
\]

对所有 \(\lambda\in\mathbf{K},\mathbf{t}\in\mathbf{T},x,y\) （在 E 中）。设 \(x\mapsto[x]\) 是 E 到由 E 的元素的形式线性组合所成的向量空间 \(\mathrm{K}_{s}^{(\mathbb{E})}\) 中的典范单射，并设 N 是 \(\mathrm{K}_{s}^{(\mathbb{E})}\) 的由元素

\[
[ \mathbf {t} + x ] - [ x ] - [ \mathbf {t} + y ] + [ y ]
\]

\[
[ \lambda \mathbf {t} + x ] - [ x ] - \lambda [ \mathbf {t} + x ] + \lambda [ x ]
\]

对 \(\lambda \in \mathbf{K}\) , \(\mathbf{t} \in \mathbf{T}\) , \(x, y\) （在 \(\mathbf{E}\) 中）；最后，令 \(\mathbf{V}\) 为 \(\mathbf{K}_s^{(\mathbf{E})}\) 模 \(\mathbf{N}\) 的商空间，并令 \(\phi\) 为 \(\mathbf{E}\) 到 \(\mathbf{V}\) 的映射，它把每个 \(x \in \mathbf{E}\) 对应到 \([x]\) 模 \(\mathbf{N}\) 的类。则 \(\phi\) 是一个仿射映射，并且对每个仿射映射 \(f\) （由 \(\mathbf{E}\) 到一个左向量空间 \(\mathbf{L}\) （在 \(\mathbf{K}\) 上）），存在且仅存在一个线性映射 \(g\) （由 \(\mathbf{V}\) 到 \(\mathbf{L}\) ）使得 \(f = g \circ \phi\) 。

\begin{enumerate}
\def\labelenumi{(\alph{enumi})}
\setcounter{enumi}{1}
\tightlist
\item
  设 \(\phi_0: \mathbf{T} \to \mathbf{V}\) 是与 \(\phi\) 相伴的线性映射，它使得 \(\phi(\mathbf{t} + x) - \phi(x) = \phi_0(\mathbf{t})\) 。证明 \(\phi\) （从而 \(\phi_0\) 也是）是单射（考虑映射 \(x \mapsto x - a\) （由 \(\mathbf{E}\) 到 \(\mathbf{T}\) ，对某个 \(a \in \mathbf{E}\) ））；对每个族 \((\lambda_t)_{t \in I}\) （ \(\mathbf{K}\) 的元素，具有有限支撑）与每个族 \((x_t)_{t \in I}\) （ \(\mathbf{E}\) 的元素），
\end{enumerate}

\[
\sum_ {i} \lambda_ {i} \phi (x _ {i}) = \phi_ {0} \left(\sum_ {i} \lambda_ {i} x _ {i}\right) \quad \text { if } \quad \sum_ {i} \lambda_ {i} = 0
\]

\[
\sum_ {i} \lambda_ {i} \phi (x _ {i}) = \mu \phi_ {0} \left(\sum_ {i} \mu^ {- 1} \lambda_ {i} x _ {i}\right) \quad \text { if } \quad \sum_ {i} \lambda_ {i} = \mu \neq 0.
\]

由此推出 \(\phi_{0}(T)\) 是 V 中过 0 的一个超平面，而 \(\phi(E)\) 是与 \(\phi_{0}(T)\) 平行的一个超平面。

\begin{enumerate}
\def\labelenumi{\arabic{enumi}.}
\setcounter{enumi}{8}
\tightlist
\item
  设 \(\mathbf{K}\) 是一个含 \(q\) 个元素的有限域，而 \(\mathbf{V}\) 是一个 \(n\) 维向量空间（在 \(\mathbf{K}\) 上）。
\end{enumerate}

\begin{enumerate}
\def\labelenumi{(\alph{enumi})}
\tightlist
\item
  证明序列 \((x_{1}, x_{2}, \ldots, x_{m})\) （ \(m \leqslant n\) 个 \(\mathbf{V}\) 的向量，构成自由系统）所成的集的基数等于
\end{enumerate}

\[
(q ^ {n} - 1) (q ^ {n} - q) \dots (q ^ {n} - q ^ {m - 1})
\]

（对 \(n\) 作归纳论证）。

\begin{enumerate}
\def\labelenumi{(\alph{enumi})}
\setcounter{enumi}{1}
\tightlist
\item
  由 (a) 推出，\(m\) 维线性簇（在一个 \(n\) 维射影空间中，空间在域 \(K\) 上）所成的集的基数等于
\end{enumerate}

\[
\frac {(q ^ {n + 1} - 1) (q ^ {n + 1} - q) \dots (q ^ {n + 1} - q ^ {m})}{(q ^ {m + 1} - 1) (q ^ {m + 1} - q) \dots (q ^ {m + 1} - q ^ {m})}.
\]

\begin{enumerate}
\def\labelenumi{\arabic{enumi}.}
\setcounter{enumi}{9}
\item
  在一个 \(n\) 维射影空间 \(\mathbf{P}(\mathbf{V})\) （在域 \(\mathbf{K}\) 上）中，射影基是一个集合 \(\mathbf{S}\) （由 \(n + 2\) 个点组成），其中任意两点都构成一个射影自由系。若 \(\mathbf{S} = (a_i)_{0 \leqslant i \leqslant n + 1}\) 和 \(\mathbf{S}' = (a_i')_{0 \leqslant i \leqslant n + 1}\) 是 \(\mathbf{P}(\mathbf{V})\) 的任意两个射影基，证明存在一个变换 \(f \in \mathbf{PGL}(\mathbf{V})\) 使得 \(f(a_i) = a_i'\) （对 \(0 \leqslant i \leqslant n + 1\) ）。为使这个变换唯一，必要且充分条件是 \(\mathbf{K}\) 可交换。（把它归结为 \(a_i' = a_i\) （对所有 \(i\) ）成立的情形，并注意总可以写成（用第 6 节的记号） \(a_i = \pi(b_i)\) ，其中 \((b_i)_{1 \leqslant i \leqslant n + 1}\) 是 \(\mathbf{V}\) 的一个基，而 \(b_0 = b_1 + b_2 + \cdots + b_{n+1}\) 。） *给出一个例子，其中 \(\mathbf{K}\) 是四元数域，并且存在点的一个无限集 \(\mathbf{T}\) （属于 \(\mathbf{P}(\mathbf{V})\) ），其中任意 \(n + 1\) 个点都构成射影自由系，又存在一个变换 \(f \in \mathbf{PGL}(\mathbf{V})\) ，它不同于恒等变换，且保持 \(\mathbf{T}\) 的所有点不变。*
\item
  设 V 是域 K 上的一个二维向量空间，而 a, b, c, d 是射影直线 \(\mathbf{P}(\mathbf{V})\) 上的四个不同点。四元组 \((a, b, c, d)\) 的交比（记作 \(\begin{bmatrix} a & b \\ d & c \end{bmatrix}\) ）是使下述条件成立的元素 \(\xi \in \mathbf{K}\) 所成的集：存在两个向量 \(u, v\) （在 \(\mathbf{V}\) 中），对它们（用第 6 节的记号）有 \(a = \pi(u)\) , \(b = \pi(v)\) , \(c = \pi(u + v)\) , \(d = \pi(u + \xi v)\) 。这个定义立即推广到具有射影直线结构的集合（第 11 节）的任意四个不同点的四元组。
\end{enumerate}

\begin{enumerate}
\def\labelenumi{(\alph{enumi})}
\item
  证明 \(\begin{bmatrix} a & b \\ d & c \end{bmatrix}\) 是一个元素 \(\neq 1\) （属于乘法群 \(\mathbf{K}^*\) ）的共轭元素的集；反之，当 \(a, b, c\) 是 \(\mathbf{P}(V)\) 的不同点而 \(\rho\) 是一个元素 \(\neq 1\) （属于 \(\mathbf{K}^*\) ）的共轭元素的集时，存在一个点 \(d \in \mathbf{P}(V)\) 使得 \(\begin{bmatrix} a & b \\ d & c \end{bmatrix} = \rho\) 。为使 \(d\) 唯一，必要且充分条件是 \(\rho\) 由单个点组成。
\item
  证明
\end{enumerate}

\[
\left[ \begin{array}{c c} a & b \\ c & d \end{array} \right] = \left[ \begin{array}{c c} b & a \\ d & c \end{array} \right] = \left[ \begin{array}{c c} a & b \\ d & c \end{array} \right] ^ {- 1}
\]

和

\[
\left[ \begin{array}{c c} d & a \\ c & b \end{array} \right] = 1 - \left[ \begin{array}{c c} a & b \\ d & c \end{array} \right]
\]

（用 \(\rho^{-1}\) （相应地， \(1 - \rho\) ）记共轭元素的集 \(\lambda \xi^{-1} \lambda^{-1} = (\lambda \xi \lambda^{-1})^{-1}\) （相应地， \(1 - \lambda \xi \lambda^{-1} = \lambda (1 - \xi) \lambda^{-1}\) ），其中 \(\xi\) 是 \(\rho\) 的一个元素）。

\begin{enumerate}
\def\labelenumi{(\alph{enumi})}
\setcounter{enumi}{2}
\tightlist
\item
  设 \((a, b, c, d)\) , \((a', b', c', d')\) 是 \(\mathbf{P}(\mathbf{V})\) 的两组由不同点组成的四元组。为使存在 V 到自身的一个双射半线性映射，使得在过渡到商时得到的 \(\mathbf{P}(\mathbf{V})\) 到自身的双射映射 f 满足条件 \(f(a) = a'\) , \(f(b) = b'\) , \(f(c) = c'\) , \(f(d) = d'\) ，必要且充分条件是存在 K 的一个自同构 \(\sigma\) 使得
\end{enumerate}

\[
\left[ \begin{array}{c c} a ^ {\prime} & b ^ {\prime} \\ d ^ {\prime} & c ^ {\prime} \end{array} \right] = \left[ \begin{array}{c c} a & b \\ d & c \end{array} \right] ^ {\sigma}.
\]

为使存在射影群 PGL(V) 的一个变换 f 满足上述条件，必要且充分条件是

\[
\left[ \begin{array}{c c} a ^ {\prime} & b ^ {\prime} \\ d ^ {\prime} & c ^ {\prime} \end{array} \right] = \left[ \begin{array}{c c} a & b \\ d & c \end{array} \right].
\]

\begin{enumerate}
\def\labelenumi{\arabic{enumi}.}
\setcounter{enumi}{11}
\tightlist
\item
  设 \(\mathbf{P}(\mathbf{V})\) 是一个有限维数 \(n\) 的（左）射影空间（在域 \(\mathbf{K}\) 上）。证明在 \(\mathbf{P}(\mathbf{V})\) 的射影超平面所成的集上存在一个 \(n\) 维（右）射影空间结构（在 \(\mathbf{K}\) 上），它典范同构于 \(\mathbf{P}(\mathbf{V}^*)\) 的结构（ \(\mathbf{V}^*\) 是 \(\mathbf{V}\) 的对偶）。若 \(\mathbf{M}\) 是维数为 \(r < n\) 的一个线性簇（在 \(\mathbf{P}(\mathbf{V})\) 中），推出一个 \((n - r - 1)\) 维射影空间结构（在包含 \(\mathbf{M}\) 的射影超平面所成的集上）。特别地，若 \(\mathbf{M}\) 的维数为 \(n - 2\) ，可以定义交比 \(\begin{bmatrix} \mathrm{H}_1 & \mathrm{H}_2 \\ \mathrm{H}_4 & \mathrm{H}_3 \end{bmatrix}\) —— 包含 M 的不同超平面的四元组 \((\mathrm{H}_1, \mathrm{H}_2, \mathrm{H}_3, \mathrm{H}_4)\) 的交比。证明若 \(\mathbf{D} \subset \mathbf{P}(\mathbf{V})\) 是一条与 M 不相交的直线，而 \(a_i\) 是 D 与 \(\mathrm{H}_i\) 的交点 ( \(1 \leqslant i \leqslant 4\) )，则
\end{enumerate}

\[
\left[ \begin{array}{c c} a _ {1} & a _ {2} \\ a _ {4} & a _ {3} \end{array} \right] = \left[ \begin{array}{c c} \mathrm{H} _ {1} & \mathrm{H} _ {2} \\ \mathrm{H} _ {4} & \mathrm{H} _ {3} \end{array} \right].
\]

*13. 在射影平面 \(\mathbf{P}(\mathbf{V})\) （在域 \(\mathbf{K}\) 上，特征 \(\neq 2\) ）中，设 \(a, b, c, d\) 是构成一个射影基的四个点（习题 10）；记 \(\mathbf{D}_{xy}\) 为过两个不同点 \(x, y\) （属于 \(\mathbf{P}(\mathbf{V})\) ）的直线，设 \(e, f, g\) 分别是直线 \(\mathbf{D}_{ab}\) 与 \(\mathbf{D}_{cd}\) 、 \(\mathbf{D}_{ac}\) 与 \(\mathbf{D}_{bd}\) 、 \(\mathbf{D}_{ad}\) 与 \(\mathbf{D}_{bc}\) 的交点；令 \(h\) 为 \(\mathbf{D}_{bc}\) 与 \(\mathbf{D}_{ef}\) 的交点；证明 \(\begin{bmatrix} b & c \\ h & g \end{bmatrix} = \{-1\}\) （“完全四边形定理”；把它归结为 \(\mathbf{D}_{ad}\) 是一个仿射平面的无穷远直线的情形）。当 \(\mathbf{K}\) 的特征为 \(2?_{*}\)

\begin{enumerate}
\def\labelenumi{\arabic{enumi}.}
\setcounter{enumi}{13}
\item
  在射影平面 \(\mathbf{P}(\mathrm{V})\) （在至少含三个元素的域 \(\mathbf{K}\) 上）中，设 \(\mathbf{D}, \mathbf{D}'\) 为两条不同的直线。为使对任意不同的点 \(a, b, c\) 属于 \(\mathbf{D}, a', b', c'\) 属于 \(\mathbf{D}'\) ，交点 \(r\) （ \(\mathbf{D}_{ab'}\) 与 \(\mathbf{D}_{ba'}\) 的交点）、 \(q\) （ \(\mathbf{D}_{ac'}\) 与 \(\mathbf{D}_{ca'}\) 的交点）、 \(p\) （ \(\mathbf{D}_{bc'}\) 与 \(\mathbf{D}_{cb'}\) 的交点）位于同一条直线上，必要且充分条件是 \(\mathbf{K}\) 可交换（“帕普斯定理”；把它归结为 \(q\) 和 \(r\) 位于一个仿射平面的无穷远直线上的情形）。把这一定理应用于 \(\mathbf{P}(\mathrm{V})\) 的直线的射影空间（习题 12）。
\item
  在域 K（至少含三个元素）上的射影平面中，设 \(s, a, b, c, a', b', c'\) 为七个不同的点，使得 \(\{s, a, b, c\}\) 和 \(\{s, a', b', c'\}\) 是射影基（习题 10），且直线 \(\mathbf{D}_{sa}, \mathbf{D}_{sb}, \mathbf{D}_{sc}\) 分别过 \(a', b', c'\) 。证明交点 \(r\) （ \(\mathbf{D}_{ab}\) 与 \(\mathbf{D}_{a'b'}\) 的交点）、 \(p\) （ \(\mathbf{D}_{bc}\) 与 \(\mathbf{D}_{b'c'}\) 的交点）、 \(q\) （ \(\mathbf{D}_{ca}\) 与 \(\mathbf{D}_{c'a'}\) 的交点）位于同一条直线上（“德萨格定理”；方法类似于习题 14）。
\item
  \begin{enumerate}
  \def\labelenumii{(\alph{enumii})}
  \tightlist
  \item
    设 \(\mathbf{E} = \mathbf{P}(\mathbf{V})\) 和 \(\mathbf{E}' = \mathbf{P}(\mathbf{V}')\) 是两个同维数 \(n \geqslant 2\) 的射影空间（分别在域 \(\mathbf{K}\) , \(\mathbf{K}'\) 上），并设 \(u\) 是 \(\mathbf{E}\) 到 \(\mathbf{E}'\) 上的一个双射映射，它把同一条直线上任意三个点映到同一条直线上的三个点。证明存在一个同构 \(\sigma\) （把 \(\mathbf{K}\) 映到 \(\mathbf{K}'\) 上）以及一个双射半线性映射 \(v\) （ \(\mathbf{V}\) 到 \(\mathbf{V}'\) 上，相对于 \(\sigma\) ），使得 \(u\) 是在过渡到商时由 \(v\) 得到的映射（“射影几何基本定理”；利用习题 7）。进一步假设 \(\mathbf{V}' = \mathbf{V}\) 且 \(\mathbf{K}\) 是交换的；为使 \(u\) 是一个射影映射，必要且充分条件是 \(\begin{bmatrix} u(a) & u(b) \\ u(d) & u(c) \end{bmatrix} = \begin{bmatrix} a & b \\ d & c \end{bmatrix}\) 对每个四元组 \((a, b, c, d)\) （由 \(\mathbf{P}(\mathbf{V})\) 中同一条直线上的不同点组成）也成立。
  \end{enumerate}
\end{enumerate}

\begin{enumerate}
\def\labelenumi{(\alph{enumi})}
\setcounter{enumi}{1}
\tightlist
\item
  设 \(p\) 是一个使 \(1 \leqslant p \leqslant n - 1\) 的整数。证明当在 \(u\) 下每个 \(p\) 维射影线性簇的像包含在一个 \(p\) 维射影线性簇中时，(a) 的第一个结论成立。
\end{enumerate}

\begin{enumerate}
\def\labelenumi{\arabic{enumi}.}
\setcounter{enumi}{16}
\tightlist
\item
  设 V 是域 K 上有限维数 n 的一个向量空间， \((e_{i})_{1\leq i\leq n}\) 是 V 的一组基， \(K'\) 是 K 的一个子域，而 \(V'\) 是 \(K'\) 上的 n 维向量空间（由 \(e_{i}\) 生成）。给出 \(V'\) 到 V 的一个单射映射的例子，它把仿射空间 \(V'\) 中同一条直线上任意三个点映到仿射空间 V 中同一条直线上的三个点，但不一定把两条平行直线映到包含在两条平行直线中的集合。（把 V 典范地嵌入与其相伴的射影空间 E 中，并考虑 E 到自身的一个射影变换 u，使得无穷远超平面在 u 下的原像不同于该超平面且不含 \(V'\) 的任何点；例如可取 K 为无限而 \(K'\) 为有限。）
\end{enumerate}

¶ 18. 设 \(\mathbf{E} = \mathbf{P}(\mathbf{V})\) 是域 \(\mathbf{K}\) 上的一个射影平面，而 \(u\) 是 \(\mathbf{E}\) 到自身的一个双射映射，它在过渡到商时由一个双射半线性映射 \(v\) （ \(\mathbf{V}\) 到自身，相对于一个自同构 \(\sigma\) —— \(\mathbf{K}\) 的自同构）导出。

\begin{enumerate}
\def\labelenumi{(\alph{enumi})}
\item
  证明以下四个性质等价：(α) 对所有 \(x \in E\) , \(x, u(x)\) 和 \(u^{2}(x)\) 位于同一条直线上；（β） E 的每条直线都含有一个在 u 下不变的点；（γ）过 E 的每个点都有一条在 u 下不变的直线；（δ）对 E 的每条直线 D，三条直线 D、 \(u(\mathbf{D})\) 和 \(u^{2}(\mathbf{D})\) 有一个公共点。（首先证明 (α) 与 (γ) 等价；利用对偶性（习题 12）推出 (β) 与 (δ) 等价；最后证明 (γ) 蕴含 (β)，并利用对偶性推出 (β) 蕴含 (γ)。）
\item
  设 u 具有 (a) 中所述的性质。证明：若在 E 中存在一条在 u 下不变的直线 D，且它仅含一个在 u 下不变的点，则 u 在过渡到商时由 V 的一个平延 v 导出（§ 10，习题 11）。（证明每条在 u 下不变的直线都含 a；通过考虑一条不过 a 的直线，证明存在一条过 a 且至少含两个在 u 下不变点的直线 \(D_{0}\) ；由此推出 \(D_{0}\) 的所有点都必然在 u 下不变。）
\item
  假设 u 具有 (a) 中所述的性质。证明：若在 E 中存在一条在 u 下不变的直线 E，且它仅含两个在 u 下不变的点，则 u 在过渡到商时由 V 的一个伸缩 v 导出（ \(\S\) 10，习题 11）。（若 a、b 是 D 中在 u 下不变的两个点，证明每条在 u 下不变的直线都过 a 或过 b；然后注意至少存在另外两个不同于 a、b 且在 u 下不变的点 c、d，因而直线 \(D_{cd}\) 过 a 或 b；通过证明 \(D_{cd}\) 的所有点都在 u 下不变来得出结论。）
\item
  假设 u 具有 (a) 中所述的性质，并且 E 中每条在 u 下不变的直线都至少含三个在 u 下不变的不同点；这时 E 中存在一个每点都在 u 下不变的射影基（习题 10）；由此推出存在 V 的一组基 \((e_{i})_{1\leq i\leq3}\) ，使得 u 在过渡到商时由一个半线性映射 \(v\) （ \(\mathbf{V}\) 到自身）导出，它满足 \(v(e_i) = e_i\) （对 \(1 \leqslant i \leqslant 3\) ）。 \(\mathbf{E}\) 中在 \(u\) 下不变的点的集这时是射影平面 \(\mathbf{P}(\mathbf{V}_0)\) ，其中 \(\mathbf{V}_0\) 是 \(\mathbf{K}_0\) （ \(\sigma\) 的不变量所成的域）上由 \(e_1, e_2, e_3\) 生成的向量空间。
\item
  此后假设 \(u\) 满足 (a) 和 (d) 的条件，且 \(u\) 和 \(u^2\) 都不是恒等映射。证明存在 \(\gamma \in \mathbf{K}\) 使得 \(\gamma^\sigma = \gamma\) 且
\end{enumerate}

\[
(\xi^ {\sigma} - \xi) ^ {\sigma} = \gamma (\xi^ {\sigma} - \xi)\tag{1}
\]

对所有 \(\xi \in \mathbf{K}\) 成立（利用条件 \((\alpha)\) （属于 \((a)\) ）以及存在 \(\zeta \in \mathbf{K}\) 使得 \(\zeta^{\sigma} \neq \zeta\) ）。证明 \(\gamma \neq -1\) ，并且对所有 \(\xi \in \mathbf{K}\) （使 \(\xi^{\sigma} \neq \xi\) ），

\[
(1 + \gamma) \xi^ {\sigma} \gamma = \gamma \xi (1 + \gamma)\tag{2}
\]

（把（1）应用于把 \(\xi\) 换成 \(\xi^2\) 的情形）；把（2）推广到所有 \(\xi \in \mathbf{K}\) ，办法是注意 \(\xi = \eta - \zeta\) （其中 \(\eta^\sigma \neq \eta\) 且 \(\xi^\sigma \neq \zeta\) ）。由此得出结论 \(\gamma \neq 1\) ，然后由（1）和（2）推出

\[
\xi^ {\sigma} = (1 + \gamma) \xi (1 + \gamma) ^ {- 1}\tag{3}
\]

对所有 \(\xi \in K\) 成立，并且 \(\gamma + \gamma^{-1}\) 属于 \(K\) 的中心。证明逆命题。*给出一个例子，其中 \(\gamma\) 不属于 \(K\) 的中心，但 \(\gamma + \gamma^{-1}\) 属于该中心。*

*19. 设 K 为一个交换域， \(\tilde{K}\) 为通过向 K 添加一个无穷远点（第 9 节）所得的射影域，而 \(f\) 和 \(g\) 为 K(X) 的两个元素。证明，记 \(h(\mathbf{X}) = f(g(\mathbf{X}))\) ，若 \(\tilde{f}, \tilde{g}, \tilde{h}\) 是 \(f, g, h\) 到 \(\tilde{\mathbf{K}}\) 的典范扩张，则 \(\tilde{h} = \tilde{f} \circ \tilde{g}_{\cdot *}\)

\BourbakiExercisesSection

\begin{enumerate}
\def\labelenumi{\arabic{enumi}.}
\tightlist
\item
  \begin{enumerate}
  \def\labelenumii{(\alph{enumii})}
  \tightlist
  \item
    设 E 是一个右 A-模。在加法群 \(E^{r}\) （ \(r \geqslant 1\) ）上定义一个外合成法则，以 \(\mathbf{M}_{r}(\mathbf{A})\) 为算子集：对每个元素 \(x = (x_{i})_{1 \leqslant i \leqslant r}\) （属于 \(E^{r}\) ）和每个方阵 \(P = (\alpha_{ij}) \in \mathbf{M}_{r}(\mathbf{A})\) ，用 x.P 记元素 \(y = (y_{i})\) （属于 \(E^{r}\) ），使得
  \end{enumerate}
\end{enumerate}

\[
y _ {i} = \sum_ {j = 1} ^ {r} x _ {j} \alpha_ {j i} \quad (1 \leqslant i \leqslant r).
\]

该外合成法则与 \(E^{r}\) 上的加法法则一起定义一个右 \(\mathbf{M}_{r}(\mathbf{A})\) -模结构（在 \(E^{r}\) 上）；把算子环限制为标量矩阵 \(I.\alpha\) 所成的环，就重新得到 \(E^{r}\) 上的积 A-模结构。证明：为使 A-模 E 有一个 r 个元素的生成系，必要且充分条件是 \(\mathbf{M}_{r}(\mathbf{A})\) -模 \(E^{r}\) 是单生成的。

\begin{enumerate}
\def\labelenumi{(\alph{enumi})}
\setcounter{enumi}{1}
\tightlist
\item
  设 \((x_{i})_{1\leqslant i\leqslant n}\) 是 E 的一个生成系，而 \((y_{j})_{1\leqslant j\leqslant m}\) 是 E 的一个元素族（相应地，E 的一个生成系）。设 \(z, z', z''\) 是 \(\mathbf{E}^{m + n}\) 的三个元素，使得 \(z_{i} = x_{i}\) 对 \(1\leqslant i\leqslant n\) , \(z_{n + j} = 0\) 对 \(1\leqslant j\leqslant m\) , \(z_{i}' = x_{i}\) 对
\end{enumerate}

\(1 \leqslant i \leqslant n, z_{n+j}^{\prime} = y_{j}\) 对 \(1 \leqslant j \leqslant n, z_{i}^{\prime \prime} = 0\) 对 \(1 \leqslant i \leqslant n, z_{n+j}^{\prime \prime} = y_{j}\) 对 \(1 \leqslant j \leqslant n\) 。证明存在两个可逆矩阵 \(P, Q\) （属于 \(\mathbf{M}_{n+m}(\mathbf{A})\) ），使得 \(z' = z.P\) （相应地， \(z'' = z.Q\) ）。

\begin{enumerate}
\def\labelenumi{(\alph{enumi})}
\setcounter{enumi}{2}
\tightlist
\item
  若 A 交换，则 \(u_{P}: x \mapsto x \cdot P\) 是 A-模 \(E^{r}\) 的一个自同态，且映射 \(P \mapsto u_{P}\) 是环 \(\mathbf{M}_{r}(\mathbf{A})\) 到环 \(\operatorname{End}_{\mathbf{A}}(\mathbf{E}^{r})\) 的一个同态。若 E 是忠实 A-模，则该同态是单射。
\end{enumerate}

\begin{enumerate}
\def\labelenumi{\arabic{enumi}.}
\setcounter{enumi}{1}
\tightlist
\item
  \begin{enumerate}
  \def\labelenumii{(\alph{enumii})}
  \tightlist
  \item
    设 \(X\) 是环 \(\mathbf{A}\) 上的一个方阵，它可以写成一个上三角矩阵 \((X_{ij})\) （由矩阵组成 \((1 \leqslant i \leqslant p, 1 \leqslant j \leqslant p)\) ）。证明：若每个方阵 \(X_{ii}\) \((1 \leqslant i \leqslant p)\) 都可逆，则 \(X\) 也可逆，且 \(X^{-1}\) 可写成一个上三角矩阵 \((Y_{ij})\) \((1 \leqslant i \leqslant p, 1 \leqslant j \leqslant p)\) ，对应于指标集的同一分划。当 \(\mathbf{A}\) 是域时，证明这一使 \(X\) 可逆的充分条件也是必要的。
  \end{enumerate}
\end{enumerate}

\begin{enumerate}
\def\labelenumi{(\alph{enumi})}
\setcounter{enumi}{1}
\tightlist
\item
  给出一个环 A 和一个矩阵 \(\begin{pmatrix} a & 1 \\ 0 & b \end{pmatrix}\) （其中 \(a \in \mathbf{A}, b \in \mathbf{A}\) ），该矩阵可逆，但 \(a\) 或 \(b\) 在 A 中均不可逆，且其逆矩阵 \(\begin{pmatrix} a' & b' \\ c' & d' \end{pmatrix}\) 满足 \(b' \neq 0\) 和 \(c' \neq 0\) （取 A 为无限维向量空间的自同态环）。
\end{enumerate}

\begin{enumerate}
\def\labelenumi{\arabic{enumi}.}
\setcounter{enumi}{2}
\tightlist
\item
  设 A 为商环 Z/30Z。证明在矩阵
\end{enumerate}

\[
\left( \begin{array}{c c c} 1 & 1 & - 1 \\ 0 & 2 & 3 \end{array} \right)
\]

上，两行线性无关，但任意两列线性相关。

\begin{enumerate}
\def\labelenumi{\arabic{enumi}.}
\setcounter{enumi}{3}
\tightlist
\item
  设 A 为一个环，C 为其中心，B 为矩阵环 \(\mathbf{M}_n(\mathbf{A}), \Delta\) 为由元素 \(\alpha\beta - \beta\alpha\) （对 \(\alpha \in \mathbf{A}, \beta \in \mathbf{A}\) ）生成的 A 的加法子群，D 为由矩阵 \(XY - YX\) （对 \(X \in \mathbf{B}, Y \in \mathbf{B}\) ）生成的 B 的加法子群； \(\Delta\) 和 D 均为 C-模。对每个矩阵 \(X = (\xi_{ij}) \in \mathbf{B}\) ，设 \(\theta(X)\) 为元素 \(\sum_{i=1}^{n} \bar{\xi}_{ij}\) （属于 A/ \(\Delta\) ），其中对所有 \(\alpha \in \mathbf{A}\) , \(\bar{\alpha}\) 表示 \(\alpha\) 模 \(\Delta\) 的类。证明 \(\theta\) 是 B 到 A/ \(\Delta\) 上的一个满同态，其核等于 D，因而 B/D 作为 C-模同构于 A/ \(\Delta\) 。（注意 D 包含矩阵 \(\alpha E_{ij}\) 以及矩阵 \(\alpha(E_{ii} - E_{jj})\) ，对 \(\alpha \in \mathbf{A}\) 和 \(i \neq j\) 。）
\end{enumerate}

¶ 5. 设 A 是一个环，L、M、N 为任意三个指标集， \(U = (a_{\lambda \mu})\) 是 \(\mathbf{A}^{\mathbf{L} \times \mathbf{M}}\) 中的一个矩阵， \(V = (b_{\mu \nu})\) 是 \(\mathbf{A}^{\mathbf{M} \times \mathbf{N}}\) 中的一个矩阵；若对每个有序对 \((\lambda, \nu) \in \mathbf{L} \times \mathbf{N}\) ，族 \((a_{\lambda \mu} b_{\mu \nu})\) （ \(\mu \in \mathbf{M}\) ）有有限支撑，则元素 \(c_{\lambda \nu} = \sum_{\mu} a_{\lambda \mu} b_{\mu \nu}\) 有定义，这时也称矩阵 \((c_{\lambda \nu})\) 是乘积 \(UV\) —— \(U\) 乘以 \(V\) 的乘积。当乘积 \(UV'\) 和 \(UV''\) 都有定义时， \(U(V' + V'')\)

和 \(UV' + UV'' = U(V' + V'')\) ；逆命题是否成立？给出三个无限矩阵 U、V、W 的例子，使得乘积 UV、VW、 \(U(VW)\) 和 \((UV)W\) 都有定义，但 \(U(VW) \neq (UV)W\) （取 U 为所有元素都等于 1 的单行矩阵，W 为其转置，V 为所有元素都是 0、1 或 -1 且每行每列只有有限个元素 \(\neq 0\) 的矩阵。用归纳法确定 V 的元素，使得 UV = 0 而 VW 只有一个元素 \(\neq 0\) 。）

\begin{enumerate}
\def\labelenumi{\arabic{enumi}.}
\setcounter{enumi}{5}
\item
  设 X 是域 K 上一个 m 行 n 列的矩阵；证明秩 \(\operatorname{rg}(X)\) 等于 X 的行数与列数相同的子矩阵的秩中的最大者。（设 \(\operatorname{rg}(X) = r\) ；若 \(a_{1}, \ldots, a_{r}\) 是 X 的 r 个列，在 \(K_{d}^{m}\) 中组成一个自由系，且用这 r 个向量和典范基的 m - r 个向量构成 \(K_{d}^{m}\) 的一组基，证明 \(a_{1}, \ldots, a_{r}\) 在典范基的其余 r 个向量上的分量构成一个秩为 r 的矩阵。）
\item
  设 X 是域 K 上的一个 m 行 n 列矩阵；若 r 是 X 的秩，证明通过删去 X 的 n - s 列得到的 m 行 s 列子矩阵的秩 \(\geqslant r + s - n\) 。
\item
  设 \(X = (\alpha_{ij})\) 是一个 \(m\) 行 \(n\) 列的矩阵（在域 \(K\) 上）。为使 \(X\) 的秩为 1，必要且充分条件是：在 \(K\) 中存在一个族 \((\lambda_i)_{1 \leqslant i \leqslant m}\) （由 \(m\) 个不全为零的元素组成）和一个族 \((\mu_j)_{1 \leqslant j \leqslant n}\) （由 \(n\) 个不全为零的元素组成），使得对每对有序指标有 \(\alpha_{ij} = \lambda_i \mu_j\) 。
\item
  设 X、Y 是域 K 上 m 行 n 列的两个矩阵；若存在两个 m 阶方阵 \(P, P_{1}\) 和两个 n 阶方阵 \(Q, Q_{1}\) ，使得 Y = PXQ 且 \(X = P_{1}XQ_{1}\) ，证明 X 与 Y 等价。
\item
  设 \(X, X', Y, Y'\) 是四个 \(n\) 阶方阵（在环 \(A\) 上），使得 \(X\) 可逆。为使存在两个可逆方阵 \(P, Q\) （ \(n\) 阶）满足 \(X' = PXQ\) 和 \(Y' = PYQ\) ，必要且充分条件是 \(X'\) 可逆且矩阵 \(YX^{-1}\) 与 \(Y'X'^{-1}\) 相似。
\item
  设 E 是域 K 上维数 \(\geqslant1\) 的一个右向量空间，H 是 E 的一个超平面。E 的每个使 H 的每个元素都不变的自同态 u 在过渡到商时给出 1 维商空间 E/H 的一个自同态，因而它具有形式 \(\dot{x}\mapsto\dot{x}\mu(\dot{x})\) ，其中 \(\mu(\dot{x})\in\mathrm{K}\) 满足 \(\mu(\dot{x}\lambda)=\lambda^{-1}\mu(\dot{x})\lambda\) （对 \(\lambda\in K^{*}\) ）。若 E 的一个使 H 的元素都不变的自同构所对应的 E/H 的自同构是恒等映射，则称它为超平面 H 的平延；否则称为超平面 H 的伸缩。当 u 是一个伸缩时，元素 \(\mu(\dot{x})\) （对 \(\dot{x}\in E/H\) ）所成的集——它是乘法群 \(K^{*}\) 中共轭元素的一个类——称为伸缩 u 的类。
\end{enumerate}

\begin{enumerate}
\def\labelenumi{(\alph{enumi})}
\item
  证明对每个伸缩，存在一条且仅一条 H 的补线，它在该伸缩下不变。
\item
  设 \(\phi\) 是 \(\mathbf{E}\) 上满足 \(\mathbf{H} = \overline{\phi}(0)\) 的一个线性形式；证明对每个平延 \(u\) （超平面 \(\mathbf{H}\) 的平延），存在唯一的向量 \(a \in \mathbf{H}\) ，使得 \(u(x) = x + a\phi(x)\) （对所有 \(x \in \mathbf{E}\) ）。设 \(\Gamma(\mathbf{E}, \mathbf{H})\) 是 \(\mathbf{GL}(\mathbf{E})\) 的由使 \(\mathbf{H}\) 的每个元素不变的自同构组成的子群；证明超平面 \(\mathbf{H}\) 的平延组成一个正规交换子群 \(\Theta(\mathbf{E}, \mathbf{H})\) （ \(\Gamma(\mathbf{E}, \mathbf{H})\) 的子群），它同构于加法群 \(\mathbf{H}\) ；商群 \(\Gamma(\mathbf{E}, \mathbf{H}) / \Theta(\mathbf{E}, \mathbf{H})\) 同构于乘法群 \(\mathbf{K}^*\) 。为使 \(\Gamma(\mathbf{E}, \mathbf{H})\) 交换，必要且充分条件是下列两个条件之一成立：（ \(\alpha\) ) \(\mathbf{K}\) 是含两个元素的域；（ \(\beta\) ) \(\mathbf{K}\) 交换且 \(\dim \mathbf{E} = 1\) 。
\item
  设 E 是有限维的；证明对每个平延 u，存在 E 的一组基，使得 u 在这组基下的矩阵的所有对角元素都等于 1，且至多还有一个别的元素 \(\neq0\) 。
\item
  证明群 \(\Theta(E, H)\) 在群 \(\mathbf{GL}(E)\) 中的中心化子（I，§ 5，第 3 号）是复合 \(Z(E)\,\Theta(E, H) = \Theta(E, H)Z(E)\) —— \(\Theta(E, H)\) 与中心 \(Z(E)\) （ \(\mathbf{GL}(E)\) 的中心）的复合（参见 § 1，习题 26）。属于该中心化子且使至少一个元素 \(\neq0\) （属于 \(E\) ）不变的自同构只有群 \(\Theta(E, H)\) 的平延。若 \(K\) 至少有 3 个元素，则 \(\Gamma(E, H)\) 在 \(\mathbf{GL}(E)\) 中的中心化子等于 \(Z(E)\) 。
\item
  证明 \(\Theta(\mathbf{E},\mathbf{H})\) 和 \(\Gamma(\mathbf{E},\mathbf{H})\) 在 \(\mathbf{GL}(\mathbf{E})\) 中的正规化子（I，§ 5，第 3 号）两者都等于由使 H 不变的自同构组成的子群。
\end{enumerate}

¶ 12. 设 E 是域 K 上维数 ≥1 的一个右向量空间。用 F(E) 记 GL(E) 的由使 1\_E - u 的核具有有限余维数的自同构 u 组成的正规子群（因而若 E 是有限维的，则 F(E) = GL(E)）。用 SL(E) 记由所有平延（习题 11）生成的 GL(E) 的正规子群；它包含在 F(E) 中。

\begin{enumerate}
\def\labelenumi{(\alph{enumi})}
\item
  若 \(\dim E \geqslant 2\) ，证明对 E 的每对非零向量 x, y，存在一个平延或两个平延的乘积，它把 x 映到 y（换句话说， \(\mathbf{SL}(\mathbf{E})\) 在 \(E = \{0\}\) 上传递地作用）。
\item
  设 V、W 是 E 的两个超平面， \(\dot{x}_{0} = x_{0} + V\) 是一个模 V 的剩余类，异于 V ，而 \(\dot{y}_{0} = y_{0} + W\) 是一个模 W 的剩余类，异于 W 。证明若 \(\dim E \geqslant 2\) ，则存在一个平延或两个平延的乘积，它把 V 映到 W 且把 \(\dot{x}_{0}\) 映到 \(\dot{y}_{0}\) （先考虑 V 与 W 相异的情形）。
\item
  若 \(\dim \mathbf{E} \geqslant 2\) ，证明任意两个异于恒等映射的平延在群 \(\mathbf{F}(\mathbf{E})\) 中共轭。
\item
  若 \(\dim \mathbf{E} \geqslant 3\) ，证明任意两个异于恒等映射的平延在群 \(\mathbf{SL}(\mathbf{E})\) 中共轭（借助 (b) 把它归结为两个平延的超平面相同的情形，然后利用 (a)）。
\item
  假设 \(\dim E = 2\) 。为使任意两个平延都在 \(\mathbf{SL}(\mathbf{E})\) 中共轭，必要且充分条件是子群 \(\mathbf{Q}\) （ \(\mathbf{K}^*\) 的、由 \(\mathbf{K}^*\) 的元素的平方生成的子群）与 \(\mathbf{K}^*\) 相同。（若 \(u\) 是一个异于恒等映射的平延， \(a\) 是 \(\mathbf{E}\) 中一个不在 \(u\) 下不变的向量，而 \(b = u(a) - a\) ，证明对每个平延 \(u'\) —— 与 \(u\) 在 \(\mathbf{SL}(\mathbf{E})\) 中共轭 —— 有 \(u'(a) - a = a\lambda + b\mu\) ，其中 \(\mu = 0\) 或 \(\mu \in \mathbf{Q}\) ；为此利用如下事实：在群 \(\mathbf{G}\) 中每个乘积 \(sts^{-1}t\) 都是平方的乘积。）
\item
  假设 \(\dim E \geqslant 2\) 。为使两个伸缩 \(u, u'\) 满足存在 \(v \in \mathbf{SL}(\mathbf{E})\) 使 \(u' = vuv^{-1}\) ，必要且充分条件是 \(u\) 与 \(u'\) 的类（习题 11）相同（利用 \((b)\) ）。由此推出，若一个伸缩的类包含在 \(\mathbf{K}^*\) 的换位子群中，则该伸缩属于群 \(\mathbf{SL}(\mathbf{E})\) （注意 \((vuv^{-1})u^{-1} = v(uv^{-1}u^{-1}))\) 。
\end{enumerate}

¶ 13. 设 E 是维数 ≥2 的右向量空间，且 H₀ 是 E 的一个超平面。

\begin{enumerate}
\def\labelenumi{(\alph{enumi})}
\item
  证明每个自同构 \(u\) （属于 \(\mathbf{F}(\mathbf{E})\) ，习题 12）是 \(\mathbf{SL}(\mathbf{E})\) 的一个自同构与可能还有超平面 \(\mathbf{H}_0\) 的一个伸缩的乘积（对 \(1_{\mathbf{E}} - u\) 的核的余维数用归纳法，并利用习题 12(a) 和 12(f)）。
\item
  证明 \(\mathbf{SL}(\mathbf{E})\) 包含 \(\mathrm{F(E)}\) 的换位子群，且除 \(\mathbf{E}\) 是含 2 个元素的域上的二维空间这一情形外与该群相同。（为证明 \(\mathbf{SL}(\mathbf{E})\) 包含 \(\mathrm{F(E)}\) 的换位子群，利用 (a) 和习题 \(12(f)\) 。为看出除所指出的情形外 \(\mathbf{SL}(\mathbf{E})\) 包含在 \(\mathrm{F(E)}\) 的换位子群中，证明对每个从 \(\mathrm{F(E)}\) 到交换群的同态，平延的像是单位元，利用习题 \(11(b)\) 和 \(12(c)\) 。）
\item
  证明 \(\mathbf{SL}(\mathbf{E})\) 等于其换位子群，除非 \(\dim \mathbf{E} = 2\) 且 \(\mathbf{K}\) 有 2 或 3 个元素。（若 \(\dim \mathbf{E} \geqslant 3\) ，注意每个平延 \(u\) 都可以写成 \(vwv^{-1}w^{-1}\) ，其中 \(v\) 是一个平延而 \(w \in \mathbf{SL}(\mathbf{E})\) ，利用习题 12(d)。若 \(\dim \mathbf{E} = 2\) ，首先注意每个自同构 \(u\) —— 它关于 \(\mathbf{E}\) 的一组基的矩阵具有形式 \(\begin{pmatrix} \lambda & 0 \\ 0 & \lambda^{-1} \end{pmatrix}\) —— 是平延的乘积，论证如 (a)；然后考虑换位子 \(uvu^{-1}v^{-1}\) ，其中 \(v\) 是平延，它关于同一组基的矩阵为 \(\begin{pmatrix} 1 & 1 \\ 0 & 1 \end{pmatrix}\) 。）
\end{enumerate}

¶ 14. 设 E 是域 K 上维数 ≥2 的向量空间。

\begin{enumerate}
\def\labelenumi{(\alph{enumi})}
\item
  证明射影群 PGL(E) 与线性群 GL(E) 对其中心（同构于 K 的中心乘法群，参见 § 1，习题 26）的商群典范同构。
\item
  设 \(\mathbf{PSL}(\mathbf{E})\) 表示群 \(\mathbf{SL}(\mathbf{E})\) 在 \(\mathbf{PGL}(\mathbf{E})\) 中的典范像；它是 \(\mathbf{PGL}(\mathbf{E})\) 的一个正规子群，包含 \(\mathbf{PGL}(\mathbf{E})\) 的换位子群（当 \(\mathbf{E}\) 有限维时）。证明 \(\mathbf{PSL}(\mathbf{E})\) 是 \(\mathbf{P}(\mathbf{E})\) 的一个双传递（参见 I，§ 5，习题 14）置换群。
\item
  设 \(a \neq 0\) 是 \(\mathbf{E}\) 的一个元素，并设 \(e = \pi(a) \in \mathbf{P}(\mathbf{E})\) （用 § 9 第 6 号的记号）。证明形如 \(x \mapsto x + a\phi(x)\) （其中 \(\phi \in \mathbf{E}^*\) 满足 \(\phi(a) = 0\) ）的平延组成 \(\mathbf{SL}(\mathbf{E})\) 的一个交换子群；若 \(\Gamma_e\) 是该子群在 \(\mathbf{PSL}(\mathbf{E})\) 中的像，则 \(\Gamma_e\) 是子群 \(\Phi_e\) （ \(\mathbf{PSL}(\mathbf{E})\) 中使 \(e\) 不变的子群）的正规子群。还证明与 \(\Gamma_e\) 共轭的子群（在 \(\mathbf{PSL}(\mathbf{E})\) 中）的并生成 \(\mathbf{PSL}(\mathbf{E})\) 。
\item
  设 \(\Sigma\) 是一个集合 F 上的本原群（I，§ 5，习题 13），它等于其换位子群。假设存在一个元素 \(c \in \mathbf{F}\) ，使得子群 \(\Phi_c\) （ \(\Sigma\) 的、使 \(c\) 不变的子群）包含一个交换正规子群 \(\Gamma_c\) ，使得与 \(\Gamma_c\) 共轭的子群（在 \(\Sigma\) 中）的并生成 \(\Sigma\) 。证明在这些条件下 \(\Sigma\) 是单群。（设 \(\Delta\) 是 \(\Sigma\) 的一个异于单位元的正规子群；利用 I，§ 5，习题 17，证明 \(\Sigma = \Delta\) . \(\Phi_c = \Delta\) . \(\Gamma_c\) 并利用 \(\Gamma_c\) 交换这一事实，推出 \(\Sigma\) 中的每个换位子都包含在 \(\Delta\) 中。）
\item
  由 (c)、(d) 和习题 13(c) 推出，群 \(\mathbf{PSL}(\mathbf{E})\) 是单群，除非 \(\dim \mathbf{E} = 2\) 且 \(\mathbf{K}\) 有 2 个或 3 个元素。
\end{enumerate}

\(^{*}(f)\) 若 \(\dim E = n + 1\) 且 K 是含 q 个元素的域，证明 PGL(E) 的阶为

\[
(q ^ {n + 1} - 1) (q ^ {n + 1} - q) \dots (q ^ {n + 1} - q ^ {n - 1}) q ^ {n}
\]

（利用(a)、§9的习题9(a)，以及K必然是交换域这一事实（V，§11，习题14和VIII，§11，第1号，定理1））。*

\begin{enumerate}
\def\labelenumi{(\alph{enumi})}
\setcounter{enumi}{6}
\item
  证明：若 \(\dim \mathbf{E} = 2\) 且 \(\mathbf{K}\) 是含 2 个（相应地，3 个）元素的域，则 \(\mathbf{PSL}(\mathbf{E})\) 同构于对称群 \(\mathfrak{S}_3\) （相应地，交错群 \(\mathfrak{A}_4\) ）。（把 \(\mathbf{PSL}(\mathbf{E})\) 看作 \(\mathbf{P}(\mathbf{E})\) 的一个置换群，注意它包含每个对换（相应地，每个循环（ \(a b c\) ）），并利用 \(f\) 。）
\item
  证明除 (g) 所考虑的情形外， \(\mathbf{SL}(\mathbf{E})\) 的每个异于 \(\mathbf{SL}(\mathbf{E})\) 的正规子群都包含在 \(\mathbf{SL}(\mathbf{E})\) 的中心中（利用 (d) 以及每个平延都是 \(\mathbf{SL}(\mathbf{E})\) 中的一个换位子（习题 13(c)））。
\end{enumerate}

¶ 15. 设 A 是一个环；对每个整数 \(n \geqslant 1\) ，把 \(\mathbf{A}^{n}\) 典范地等同于 \(\mathbf{A}^{(\mathbf{N})}\) 的由指标 \(\geqslant n\) 的坐标为 0 的元素组成的子模，并用 \(\mathbf{A}_{n}^{\prime}\) 记补子模，它由指标 \(<n\) 的坐标为 0 的元素组成。每个自同态 \(u \in \mathbf{GL}_{n}(\mathbf{A})\) 等同于 \(\mathbf{A}^{(\mathbf{N})}\) 的这样一个自同构：它在 \(\mathbf{A}^{n}\) 上的限制是 \(u\) ，在 \(\mathbf{A}_{n}^{\prime}\) 上的限制是恒等映射，于是 \(\mathbf{GL}_{n}(\mathbf{A})\) 被等同于 \(\mathbf{GL}(\mathbf{A}^{(\mathbf{N})})\) 的一个子群；用 F 记 \(\mathbf{GL}(\mathbf{A}^{(\mathbf{N})})\) 的由诸 \(\mathbf{GL}_{n}(\mathbf{A})\) 的并集组成的子群，换言之，即除有限个外使 \(\mathbf{A}^{(\mathbf{N})}\) 的典范基的元素都不变的子群（参见习题 12）。

设 \(\mathbf{T}_n\) 是 \(\mathbf{GL}_n(\mathbf{A})\) 的由诸 \(\mathrm{B}_{ij}(\lambda)\) （对 \(0 \leqslant i < n\) , \(0 \leqslant j < n\) , \(i \neq j\) , \(\lambda \in \mathbf{A}\) ，第 13 号）生成的子群，并令 \(\mathbf{T}\) 为 \(\mathbf{F}\) 的由诸 \(\mathbf{T}_n\) 的并集组成的子群。

\begin{enumerate}
\def\labelenumi{(\alph{enumi})}
\item
  证明对 \(n \geqslant 3\) ， \(T_n\) 等于其换位子群。
\item
  设 \(a, b\) 是 \(A\) 的两个可逆元素。证明
\end{enumerate}

\[
\left( \begin{array}{c c} a b & 0 \\ 0 & 1 \end{array} \right) = P \left( \begin{array}{c c} a & 0 \\ 0 & b \end{array} \right) Q \quad \text { and } \quad \left( \begin{array}{c c} a & 0 \\ 0 & b \end{array} \right) = P ^ {\prime} \left( \begin{array}{c c} b & 0 \\ 0 & a \end{array} \right) Q ^ {\prime}
\]

其中 \(P, Q, P', Q'\) 属于 \(T_2\) 。由此推出，矩阵 \(\begin{pmatrix} aba^{-1}b^{-1} & 0 \\ 0 & 1\end{pmatrix}\) 属于 \(T_2\) 。

\begin{enumerate}
\def\labelenumi{(\alph{enumi})}
\setcounter{enumi}{2}
\tightlist
\item
  由 (b) 推出， \(\mathbf{GL}_{n}(\mathbf{A})\) 的换位子群包含在 \(T_{2n}\) 中（在 (b) 中把 A 换成 \(\mathbf{M}_{n}(\mathbf{A})\) ）。（利用 (a)）得出结论：F 的换位子群等于 T。
\end{enumerate}

\begin{enumerate}
\def\labelenumi{\arabic{enumi}.}
\setcounter{enumi}{15}
\tightlist
\item
  \begin{enumerate}
  \def\labelenumii{(\alph{enumii})}
  \tightlist
  \item
    设 \(U\) 是一个 \(n\) 阶方阵（在域 \(\mathbf{Q}\) 上），其所有非对角元素都等于同一个有理数 \(r > 0\) ，且其对角线 \((d_1, \ldots, d_n)\) 满足： \(d_i \geqslant r\) （对所有 \(i\) ）且 \(d_i > r\) （除 \(i\) 的可能一个值外）。证明 \(U\) 可逆（若 \(x = (x_i)_{1 \leqslant i \leqslant n}\) 是方程 \(U.x = 0\) 的一个解，证明这些 \(x_i\) 都有相同的符号，并由此推出它们全为零）。
  \end{enumerate}
\end{enumerate}

\begin{enumerate}
\def\labelenumi{(\alph{enumi})}
\setcounter{enumi}{1}
\tightlist
\item
  设 E 是含 m 个元素的有限集，这些元素记作 \(a_{i}\) ( \(1 \leqslant i \leqslant m\) )，而 \((\mathbf{A}_{j})_{1 \leqslant j \leqslant n}\) 是 E 的 n 个子集所成的族，彼此互不相同；假设 \(\text{Card}(\mathbf{A}_{i} \cap \mathbf{A}_{j}) = r \geqslant 1\) 对每个有序对 \((i, j)\) （由 \([1, n]\) 的不同元素组成的有序对）成立。证明必然有 \(n \leqslant m\) 。（令 \(V = (c_{ij})\) 是 \((m, n)\) 型矩阵，使得 \(c_{ij} = 1\) （当 \(a_{i} \in A_{j}\) 时），否则 \(c_{ij} = 0\) 。把 (a) 应用于矩阵 \(^{t}V\) . V （n 阶）。）
\end{enumerate}

¶ 17. 设 K 为一个域（交换或非交换），E 为 K 上的 n 维右向量空间，u 为 E 的一个自同态。

\begin{enumerate}
\def\labelenumi{(\alph{enumi})}
\item
  证明：若 \(n > 1\) 且 \(u\) 不是中心位似，则存在 \(E\) 的一组基，使得 \(u\) 关于这组基的矩阵具有形式 \((\alpha_{ij})\) ，其中 \(\alpha_{ij} = 0\) （对 \(j > i + 1\) ）且 \(\sum_{i=1}^{n-1} \alpha_{i,i+1} = 1\) （对 \(n\) 用归纳法论证）。
\item
  证明：若 \(n > 2\) ，则存在 \(\mathbf{E}\) 的一组基，使得 \(u\) 关于这组基的矩阵具有形式 \((\alpha_{ij})\) ，其中 \(\alpha_{ij} = 0\) （对 \(j > i + 1\) ）且
\end{enumerate}

\(\sum_{i=1}^{n-1} \alpha_{i, i+1} = 0\) 。（对 \(n\) 用归纳法并利用 \((a)\) 论证，注意对 \(u\) 或对 \(u + \gamma 1_{E}\) （其中 \(\gamma\) 属于 \(K\) 的中心）证明该命题是等价的，并分别研究 \(u\) 的秩为 1 或 2 的情形。）

\begin{enumerate}
\def\labelenumi{(\alph{enumi})}
\setcounter{enumi}{2}
\item
  设 \(A = (\alpha_{ij})\) 是具有 (b) 中所述性质的一个矩阵。又设 \(S\) 是阶为 \(n(\varepsilon_{ij})\) 的矩阵，使得 \(\varepsilon_{ij} = 0\) ，除非 \(i = j + 1\) ，这时 \(\varepsilon_{j+1,j} = 1\) 。证明存在一个矩阵 \(B = (\beta_{ij})\) （ \(n\) 阶），使得 \(A = \lambda E_{nn} + (BS - SB)\) （ \(\lambda \in \mathbf{K}\) ）。（取 \(B\) 使得 \(\beta_{i1} = 0\) （对 \(1 \leqslant i \leqslant n\) ）且 \(\beta_{ij} = 0\) （对 \(j > i + 2\) ）。）
\item
  假设 K 是交换的。由 (c) 推出，E 的每个满足 \(\operatorname{Tr}(u) = 0\) 的自同态 u 都可以写成 vw - wv，其中 v 和 w 是 E 的自同态。
\end{enumerate}

\BourbakiExercisesSection

\begin{enumerate}
\def\labelenumi{\arabic{enumi}.}
\tightlist
\item
  设 \(\Delta\) 是一个交换幺半群，用加法记号书写，其单位元记为 0。设 \(\mathbf{Z}^{(\Delta)}\) 是 \(\Delta\) 在 \(\mathbf{Z}\) 上的代数（III，§2，第 6 号），即一个交换环，它是自由 \(\mathbf{Z}\) -模，具有基 \((\mathrm{X}^{\sigma})_{\sigma \in \Delta}\) 和乘法 \((a\mathrm{X}^{\sigma})(b\mathrm{X}^{\tau}) = ab\mathrm{X}^{\sigma +\tau}\) 。对每个 \(\mathbf{Z}\) -模，我们记 \(\mathrm{E}^{(\Delta)} = \mathrm{E} \otimes_{\mathbf{Z}} \mathbf{Z}^{(\Delta)}\) ，因而 \(\mathrm{E}^{(\Delta)}\) 的每个元素都可唯一地写成形式 \(\sum_{\sigma \in \Delta} z_{\sigma} \otimes \mathrm{X}^{\sigma}\) ，其中 \(z_{\sigma} \in \mathrm{E}\) ，族 \((z_{\sigma})\) 有有限支撑。对每个 \(\mathbf{Z}\) -模同态 \(f: \mathrm{E} \to \mathrm{E}'\) ，用 \(f^{(\Delta)}\) 记 \(\mathbf{Z}^{(\Delta)}\) -模同态 \(f \otimes 1: \mathrm{E}^{(\Delta)} \to \mathrm{E}'^{(\Delta)}\) 。类似地，若 \(\Delta'\) 是另一个单位元记为 0 的加法幺半群，且 \(\alpha: \Delta \to \Delta'\) 是一个幺半群同态，则用 \(\alpha(\mathrm{E})\) 记同态 \(\mathrm{E}^{(\Delta)} \to \mathrm{E}^{(\Delta')}\) ，它满足 \(\alpha(\mathrm{E})(z_{\sigma} \otimes \mathrm{X}^{\sigma}) = z_{\sigma} \otimes \mathrm{X}^{\alpha(\sigma)}\) 。把 \((\mathrm{E}^{(\Delta)})^{(\Delta')}\) 与 \(\mathrm{E}^{(\Delta \times \Delta')}\) 等同，办法是把 \(\sum_{\sigma' \in \Delta'} \left( \sum_{\sigma \in \Delta} z_{\sigma, \sigma'} \otimes \mathrm{X}^{\sigma} \right) \otimes \mathrm{X}^{\sigma'}\) 对应到元素 \(\sum_{\sigma, \sigma'} z_{\sigma, \sigma'} \otimes \mathrm{X}^{(\sigma, \sigma')}.\) 最后，用 \(\varepsilon(\mathrm{E})\) 记同态 \(\mathrm{E}^{(\Delta)} \to \mathrm{E}\) ，它把 \(\sum_{\sigma} z_{\sigma} \otimes \mathrm{X}^{\sigma}\) 映到 \(\sum_{\sigma} z_{\sigma}\)
\end{enumerate}

\begin{enumerate}
\def\labelenumi{(\alph{enumi})}
\tightlist
\item
  设 E 是一个 Z-模， \((\mathrm{E}_{\sigma})_{\sigma\in\Delta}\) 是 E 上的一个 \(\Delta\) 型分次，且对所有 \(\sigma\in\Delta\) ，设 \(p_{\sigma}\) 是投影算子 \(E\to E_{\sigma}\) ，它对应于把 E 分解为诸 \(E_{\sigma}\) 的直和。对所有 \(x\in E\) ，设 \(\phi_{\mathbf{E}}(x)=\sum_{\sigma\in\Delta}p_{\sigma}(x)\otimes X^{\sigma}\) 。证明这样定义的同态 \(\phi_{E}\colon E\to E^{(\Delta)}\) 具有下列性质：
\end{enumerate}

\[
\varepsilon (\mathbf {E}) \circ \phi_ {\mathbf {E}} = 1 _ {\mathbf {E}}\tag{1}
\]

\[
\delta (\mathbf {E}) \circ \phi_ {\mathbf {E}} = \phi_ {\mathbf {E}} ^ {(\Delta)} \circ \phi_ {\mathbf {E}},\tag{2}
\]

其中 \(\delta: \Delta \to \Delta \times \Delta\) 是对角映射。

证明这样就定义了 E 上的 \(\Delta\) 型分次所成的集到 E 到 \(\mathrm{E}^{(\Delta)}\) 中的满足条件 (1) 和 (2) 的同态所成的集上的一个双射。

\begin{enumerate}
\def\labelenumi{(\alph{enumi})}
\setcounter{enumi}{1}
\item
  设 E、F 是两个 \(\Delta\) 型分次 Z-模。为使同态 \(f: E \to F\) 是 0 次分次的，必要且充分条件是 \(\phi_{F} \circ f = f^{(\Delta)} \circ \phi_{E}\) 。为使 E 的子模 \(E'\) 是分次的，必要且充分条件是 \(\phi_{E}(E') \subset E'^{(\Delta)}\) 。
\item
  设 \(\mathbf{A}\) 是一个环；给 \(\mathbf{A}^{(\Delta)} = \mathbf{A} \otimes_{\mathbf{Z}} \mathbf{Z}^{(\Delta)}\) 赋予一个环结构，使得 \((a \otimes \mathbf{X}^{\sigma})(b \otimes \mathbf{X}^{\tau}) = (ab) \otimes \mathbf{X}^{\sigma + \tau}\) 。设 \((\mathbf{A}_{\sigma})_{\sigma \in \Delta}\) 是加法群 A 的一个分次；对于这个分次与 A 上的环结构以及
\end{enumerate}

\(1 \in A_{0}\) 相容，必要且充分条件是 \(\phi_{A}: A \to A^{(\Delta)}\) 是一个环同态。陈述并证明关于分次 A-模的一个类似结果。

\BourbakiAppendixExercises

\begin{enumerate}
\def\labelenumi{\arabic{enumi}.}
\tightlist
\item
  设 A 是一个伪环，M 是 A 上的一个左伪模，使得对所有 \(z \neq 0\) （在 M 中），有 \(z \in Az\) 。证明若 x, y 是 M 的两个元素，使得在 A 中 \(\text{Ann}(x) \subset \text{Ann}(y)\) ，则存在 M 的一个自同态 u，使得 \(u(x) = y\) （证明关系 bx = ax（对 A 中的 a, b）蕴含 by = ay）。
\end{enumerate}

\chapter{张量代数、外代数、对称代数}

回顾第一章中引入的指数记号，我们将频繁使用它（I，§ 7，第 8 号）：

设 \((x_{\lambda})_{\lambda \in L}\) 是环 \(\mathbf{A}\) 的一族两两可交换的元素；对每个具有有限支撑的映射 \(\alpha: \mathbf{L} \to \mathbf{N}\) ，我们将记

\[
x ^ {\alpha} = \prod_ {\lambda \in L} x _ {\lambda} ^ {\alpha (\lambda)}.
\]

若 \(\beta\) 是另一个 \(\mathbf{L}\) 到 \(\mathbf{N}\) 的具有有限支撑的映射，则 \(\alpha + \beta\) 表示映射

\[
\lambda \mapsto \alpha (\lambda) + \beta (\lambda)
\]

把 L 映入 N；带有这个合成法则，由 L 到 N 的具有有限支撑的映射所成的集 \(\mathbf{N}^{(L)}\) 是由 L 导出的自由交换幺半群，而

\[
x ^ {\alpha} x ^ {\beta} = x ^ {\alpha + \beta},
\]

对所有 \(\alpha \in \mathbf{N}^{(L)}\) ，我们记 \(|\alpha| = \sum_{\lambda \in L} \alpha(\lambda) \in \mathbf{N}\) ；于是 \(|\alpha + \beta| = |\alpha| + |\beta|\) ； \(|\alpha|\) 称为“多重指标” \(\alpha\) 的阶。对所有 \(\lambda \in L\) ，用 \(\delta_{\lambda}\) 记 \(\mathbf{N}^{(L)}\) 中使得 \(\delta_{\lambda}(\lambda) = 1\) , \(\delta_{\lambda}(\mu) = 0\) （对 \(\mu \neq \lambda\) ）的元素（克罗内克指标）；诸 \(\delta_{\lambda}\) （对 \(\lambda \in L\) ）是 \(\mathbf{N}^{(L)}\) 中仅有的阶为 1 的元素。给 \(\mathbf{N}^{(L)}\) 赋予由 \(\mathbf{N}^{\mathrm{L}}\) 上的乘积序诱导的序，使得关系 \(\alpha \leqslant \beta\) 等价于“ \(\alpha(\lambda) \leqslant \beta(\lambda)\) （对所有 \(\lambda \in L\) ）”；这时多重指标 \(\lambda \mapsto \beta(\lambda) - \alpha(\lambda)\) 记作 \(\beta - \alpha\) ，因而它是使 \(\alpha + (\beta - \alpha) = \beta\) 的唯一多重指标。对所有 \(\alpha \in \mathbf{N}^{(L)}\) ，只有有限多个多重指标 \(\beta \leqslant \alpha\) ；诸 \(\delta_{\lambda}\) 是集合 \(\mathbf{N}^{(L)} = \{0\}\) 的极小元素；关系 \(\alpha \leqslant \beta\) 蕴含 \(|\alpha| \leqslant |\beta|\) ，且若 \(\alpha \leqslant \beta\) 与 \(|\alpha| = |\beta|\) 两者都成立，则 \(\alpha = \beta\) 。最后，我们记 \(\alpha! = \prod_{\lambda \in L} (\alpha(\lambda))!\) ，它有意义，因为 \(0! = 1\) 。

从第4节到第8节（含），A表示一个交换环，且除非另有说明，所考虑的代数均假定为结合的且有单位的，代数同态也假定为有单位的。

\section{代数}
